\documentclass[12pt,a4paper,openany]{report}

\usepackage[utf8]{inputenc}
\usepackage[T1]{fontenc}
\usepackage[french]{babel}
\usepackage{csquotes}
\usepackage{libertinus}
\usepackage{microtype}
\usepackage{amsmath,amssymb,mathtools}
\usepackage{booktabs,tabularx,array,colortbl}
\usepackage{graphicx}
\usepackage{placeins}
\usepackage{tikz}
\usetikzlibrary{arrows.meta,positioning,calc}
\usepackage{pgfplots}
\pgfplotsset{compat=1.18}
\pgfplotsset{courseplot/.style={width=.82\textwidth,height=7.2cm,
  grid=major,grid style={gridgray!40}}}
\usepackage{xcolor}
\usepackage{enumitem}
\usepackage{tcolorbox}
\tcbuselibrary{breakable,skins}
\usepackage[margin=2.45cm]{geometry}
\usepackage[backend=biber,style=numeric,sorting=none,defernumbers=true]{biblatex}
\addbibresource{references.bib}
\AtEveryBibitem{\clearfield{urldate}}
\DeclareBibliographyCategory{reading}
\addtocategory{reading}{bishop2006,ng2018cs229,murphy2022,
  james2023,openintro2024,stanfordcs109,psustat504}
\usepackage{hyperref}

\definecolor{ink}{HTML}{18212B}
\definecolor{accent}{HTML}{155E75}
\definecolor{accentlight}{HTML}{E7F3F6}
\definecolor{warning}{HTML}{9A3412}
\definecolor{warninglight}{HTML}{FFF1E8}
\definecolor{gridgray}{HTML}{CBD5E1}
\hypersetup{colorlinks=true,linkcolor=ink,citecolor=accent,urlcolor=accent,
  pdftitle={Régression logistique et classification un-contre-tous}}
\setlist{nosep,leftmargin=*}
\setlength{\parindent}{0pt}
\linespread{1.08}
\setlength{\parskip}{0.72em}
\setlength{\abovedisplayskip}{0.9em}
\setlength{\belowdisplayskip}{0.9em}
\setlength{\abovedisplayshortskip}{0.65em}
\setlength{\belowdisplayshortskip}{0.65em}
\setlength{\textfloatsep}{1.25em plus 0.25em minus 0.2em}
\setlength{\floatsep}{1.1em plus 0.2em minus 0.2em}
\renewcommand{\arraystretch}{1.16}
\setcounter{tocdepth}{2}
% Un encadré reporté laisse du blanc en fin de page : mieux vaut ce blanc qu'un
% étirement des interlignes sur toute la page.
\raggedbottom
\widowpenalty=10000
\clubpenalty=10000

% Aucun de ces encadrés ne dépasse une page : les laisser sécables produit des
% blocs coupés après deux lignes. Seul workflow, liste de dix étapes, reste
% breakable.
\newtcolorbox{resultat}[1][]{enhanced,colback=accentlight,
  colframe=accent,boxrule=0.7pt,arc=1mm,left=2mm,right=2mm,top=1.5mm,
  bottom=1.5mm,title={#1},fonttitle=\bfseries,coltitle=white}
\newtcolorbox{vigilance}[1][]{enhanced,colback=warninglight,
  colframe=warning,boxrule=0.7pt,arc=1mm,left=2mm,right=2mm,top=1.5mm,
  bottom=1.5mm,title={#1},fonttitle=\bfseries,coltitle=white}
\newtcolorbox{lecture}{enhanced,colback=gridgray!12,
  colframe=ink!55,boxrule=0.55pt,arc=1mm,left=2mm,right=2mm,top=1.5mm,
  bottom=1.5mm}
\newtcolorbox{procedure}[1]{enhanced,colback=white,
  colframe=accent!75,boxrule=0.65pt,arc=1mm,left=2mm,right=2mm,top=1.5mm,
  bottom=1.5mm,title={Sous-procédure --- #1},fonttitle=\bfseries,coltitle=white}
\newtcolorbox{pseudo}[1][]{enhanced,colback=white,
  colframe=ink!60,boxrule=0.6pt,arc=1mm,left=2mm,right=2mm,top=1.5mm,
  bottom=1.5mm,title={#1},fonttitle=\bfseries,coltitle=white,
  fontupper=\ttfamily\small,before upper=\setlength{\parskip}{0pt}}
\newtcolorbox{fichier}[1][]{enhanced,colback=ink!5,
  colframe=ink!45,boxrule=0.55pt,arc=1mm,left=2mm,right=2mm,top=1.5mm,
  bottom=1.5mm,title={#1},fonttitle=\bfseries,coltitle=white,
  fontupper=\ttfamily\small,before upper=\setlength{\parskip}{0pt}}
\newtcolorbox{workflow}[1]{enhanced,breakable,colback=accentlight,
  colframe=accent,boxrule=0.9pt,arc=1mm,left=2mm,right=2mm,top=1.5mm,
  bottom=1.5mm,title={#1},fonttitle=\bfseries,coltitle=white}

\newcommand{\dsym}[2]{\hypertarget{sym:#1}{#2}}
\newcommand{\rsym}[2]{\hyperlink{sym:#1}{\textcolor{accent}{#2}}}
% Même lien, dans un bandeau de titre : la couleur accent y serait illisible.
\newcommand{\rsymt}[2]{\hyperlink{sym:#1}{\textcolor{white}{#2}}}

\DeclareMathOperator*{\Argmax}{arg\,max}
\newcommand{\R}{\mathbb{R}}
\newcommand{\Pp}{\mathbb{P}}
\newcommand{\ind}{\mathbf{1}}
\newcommand{\x}{\mathbf{x}}
\newcommand{\xt}{\tilde{\mathbf{x}}}
\newcommand{\X}{\mathbf{X}}
\newcommand{\y}{\mathbf{y}}
\newcommand{\thv}{\boldsymbol{\theta}}
\newcommand{\eps}{\varepsilon}

\title{\textbf{Régression logistique et classification un-contre-tous}\\
  \large Modélisation, estimation, calcul numérique et évaluation}
\author{Sylvain Maitre}
\date{\today}

\begin{document}
\hypersetup{pageanchor=false}
\maketitle
\pagenumbering{roman}

\begin{abstract}
Un fichier contient mille six cents élèves. De chacun, on connaît treize notes et
la maison. Il faut en tirer une règle qui donne sa maison à un élève dont on ne
connaît que les notes.

La règle construite ici n'annonce pas directement une maison : elle calcule une
probabilité. Cette probabilité vient d'une somme pondérée des notes, ramenée
entre zéro et un par la fonction logistique. Les poids de cette somme sont les
inconnues. Le principe du maximum de vraisemblance les fixe : on retient les
poids sous lesquels les maisons effectivement observées deviennent les plus
probables. Cette exigence s'écrit comme un coût à réduire, dont la pente indique
en chaque point la correction à appliquer aux poids ; répéter la correction fait
descendre le coût jusqu'à son minimum.

Quatre maisons demandent quatre exécutions de ce procédé, chacune posant la
question « celle-ci, ou bien l'une des trois autres ». Un élève reçoit alors
quatre scores et se voit attribuer la maison du plus grand ; softmax n'est pas
nécessaire pour décider. Le reste du travail porte sur les mesures : quelles
notes retenir, quoi mettre dans les cases vides, à quoi reconnaître un résultat
qui vaudra encore pour des élèves jamais vus. Cette part occupe autant de place
que le modèle lui-même, un classifieur qui reconnaît sans faute les élèves de
son propre fichier n'ayant encore rien démontré.

Ce texte a été rédigé avec GPT et Claude selon mes instructions et éléments fournis. Le résultat n'est pas formidable mais il est lisible et peut apporter des éléments de compréhension. Il est mis à disposition sous licence CC-BY-SA 4.0.
\end{abstract}

\tableofcontents
\clearpage
\pagenumbering{arabic}
\hypersetup{pageanchor=true}
\chapter{Problème statistique et chaîne de traitement}
\label{ch:problem}

\begin{lecture}
Le fichier fournit des élèves dont la maison est connue. Le modèle cherche des
régularités reliant leurs notes à leur maison, puis les applique à un élève
absent du fichier. Toute l'analyse porte sur l'écart entre la réussite sur les
élèves déjà vus et la réussite sur de nouveaux élèves.
\end{lecture}

On observe $m$ élèves. L'élève d'indice $i$ est décrit par $n$ variables numériques
$\x_i\in\R^n$ et par une maison $c_i$ appartenant à
\[
  \mathcal C=\{\text{Gryffindor},\text{Hufflepuff},
  \text{Ravenclaw},\text{Slytherin}\},\qquad K=|\mathcal C|=4.
\]
Le but est une règle $f:\R^n\to\mathcal C$ dont l'erreur reste faible sur des
élèves nouveaux, issus du même mécanisme de génération des données. La
performance sur l'échantillon observé sert d'indicateur intermédiaire.

Le terme \emph{régression} désigne ici la modélisation d'un log-rapport de chances
par une fonction affine. La variable réponse reste qualitative, ce qui range le
problème dans la classification \cite[chap.~4]{hastie2009}.

La table~\ref{tab:sample} montre quatre lignes du fichier. Ces valeurs passent par
l'imputation et la standardisation du chapitre~\ref{ch:data} avant d'entrer dans
le modèle. La colonne House porte la réponse à prédire.

\begin{table}[htbp]
\centering
\begin{tabular}{lrrrrl}
\toprule
élève & Astronomy & Herbology & Charms & Flying & House \\
\midrule
$1$ & $-487{,}9$ & $ 5{,}7$ & $-232{,}8$ & $-26{,}9$  & Ravenclaw \\
$2$ & $-552{,}1$ & $-6{,}0$ & $-252{,}2$ & $-113{,}5$ & Slytherin \\
$3$ & $-366{,}1$ & $ 7{,}7$ & $-227{,}3$ & $ 30{,}4$  & Ravenclaw \\
$4$ & $ 697{,}7$ & $-6{,}5$ & $-256{,}8$ & $200{,}6$  & Gryffindor \\
\bottomrule
\end{tabular}
\caption{Quatre lignes du jeu d'entraînement, réduites à quatre colonnes de notes
sur les treize disponibles. Une ligne représente une observation.}
\label{tab:sample}
\end{table}

\begin{figure}[htbp]
  \centering
  \resizebox{\textwidth}{!}{%
  \begin{tikzpicture}[
    node distance=5mm and 5mm,
    block/.style={draw=accent,rounded corners=1mm,fill=accentlight,
      align=center,minimum height=10mm,text width=27mm},
    arr/.style={-{Latex[length=2mm]},thick,draw=ink}]
    \node[block] (raw) {données\\brutes};
    \node[block,right=of raw] (split) {partition\\train--test};
    \node[block,right=of split] (prep) {imputation et\\standardisation};
    \node[block,right=of prep] (fit) {$K$ modèles\\binaires};
    \node[block,right=of fit] (eval) {décision et\\évaluation};
    \draw[arr] (raw)--(split); \draw[arr] (split)--(prep);
    \draw[arr] (prep)--(fit); \draw[arr] (fit)--(eval);
  \end{tikzpicture}}
  \caption{Chaîne de traitement. Les statistiques de préparation sont estimées sur
  l'ensemble d'apprentissage, puis appliquées telles quelles aux données de test.}
\label{fig:pipeline}
\end{figure}

\begin{workflow}{Ordre global : apprentissage, sélection et évaluation}
\begin{enumerate}
  \item \textbf{Partitionner.} Construire les indices d'apprentissage et de test,
        puis mettre le test à l'écart.
  \item \textbf{Préparer.} Estimer imputation, moyennes et écarts-types sur
        l'apprentissage ; transformer les données avec ces valeurs.
  \item \textbf{Organiser la sélection.} Réserver une validation ou construire
        des plis à l'intérieur des données disponibles pour l'apprentissage.
  \item \textbf{Former les cibles OvR.} Pour chaque classe $k$, construire le
        vecteur binaire $\y_k$.
  \item \textbf{Ajuster une configuration.} Pour chaque classe, initialiser
        $\thv_k$, puis répéter scores, probabilités, gradient et mise à jour
        jusqu'au critère d'arrêt.
  \item \textbf{Comparer les configurations.} Choisir $\alpha$, $\lambda$, le
        seuil et les autres hyperparamètres sur la seule validation.
  \item \textbf{Figer le modèle.} Conserver la préparation, l'ordre des
        caractéristiques, les $K$ vecteurs $\thv_k$ et l'ordre des classes.
  \item \textbf{Prédire le test.} Appliquer la préparation conservée, calculer
        les $K$ scores de chaque ligne et retenir leur maximum.
  \item \textbf{Évaluer une fois.} Comparer prédictions et classes réelles du
        test ; produire matrice de confusion et mesures par classe.
  \item \textbf{Archiver l'expérience.} Conserver graine, indices,
        hyperparamètres, historiques de convergence et résultats finaux.
\end{enumerate}
\end{workflow}

Chaque cadre « sous-procédure » du document développe une opération appelée par
cet ordre global : un calcul interne, un contrôle ou une méthode de sélection.

\begin{resultat}[Risque de généralisation]
La quantité pertinente est le \rsym{risque}{risque de généralisation}
$R(f)=\Pp(f(X)\neq C)$. La loi jointe de $(X,C)$ étant inconnue, on l'estime sur
un ensemble de test tenu à l'écart du choix des variables, de l'estimation des
transformations et de l'ajustement des paramètres.
\end{resultat}

Sur un test indépendant de $200$ élèves dont $160$ sont correctement classés,
l'erreur observée vaut $40/200=0{,}20$. Elle estime le risque de généralisation :
environ $20\,\%$ des futurs élèves comparables seraient mal classés. L'estimation
porte sa propre variabilité, un autre échantillon de $200$ élèves donnant un
résultat voisin et différent.

\section{Décision et calibration probabiliste}

Attribuer une maison à un nouvel élève est une décision. Estimer une probabilité
d'appartenance \rsym{calibration}{calibrée} est une tâche plus exigeante : parmi les élèves auxquels
le modèle attribue une probabilité proche de $0{,}8$, environ $80\,\%$
appartiennent effectivement à la classe. OvR répond à la première question ; la
seconde demande les précautions du chapitre~\ref{ch:ovr}.

Supposons que le modèle attribue une probabilité comprise entre $0{,}75$ et
$0{,}85$ à Ravenclaw pour $50$ élèves. Si $39$ sont réellement Ravenclaw, la
fréquence observée vaut $39/50=0{,}78$ et confirme l'annonce. Si $20$ le sont,
elle tombe à $0{,}40$ : le classement par maximum des scores reste utilisable, et
la lecture probabiliste de ces nombres demande une recalibration préalable.

\begin{resultat}[Objets définis par le problème]
Le jeu étiqueté fournit une matrice de caractéristiques et un vecteur de maisons.
L'apprentissage produit des paramètres. La prédiction transforme une ligne de
caractéristiques en une maison. L'évaluation compare le vecteur des maisons
prédites au vecteur des maisons réelles. Ces objets gardent des rôles séparés
dans toute la chaîne.
\end{resultat}

\chapter{Données, préparation et protocole expérimental}
\label{ch:data}

\section{Matrice de design}

\begin{lecture}
Une matrice de design est le tableau numérique reçu par le modèle. Chaque ligne
décrit un élève, chaque colonne une caractéristique. Une colonne supplémentaire
remplie de $1$ ajoute le même terme constant au score de tous les élèves.
\end{lecture}

Après préparation, on ajoute une coordonnée constante à chaque observation :
\[
  \xt_i=(1,x_{i1},\ldots,x_{in})^\top\in\R^{n+1}.
\]
La matrice de design empile les observations par lignes,
\[
  \X=\begin{bmatrix}\xt_1^\top\\ \vdots\\\xt_m^\top\end{bmatrix}
  \in\R^{m\times(n+1)}.
\]
Cette convention fixe toutes les dimensions ultérieures : les paramètres sont
des vecteurs colonnes de $\R^{n+1}$ et $\X\thv\in\R^m$.

Sur trois élèves décrits par deux notes standardisées, la construction devient
par exemple
\begin{figure}[htbp]
\centering
\begin{tikzpicture}[
  box/.style={draw=ink!55,rounded corners=2pt,fill=accentlight,
    inner xsep=7pt,inner ysep=5pt,align=center}]
  \node[box] (raw) {$
    \begin{array}{c|rr}
      &x_1&x_2\\\hline
      \text{élève 1}&-1{,}2&0{,}4\\
      \text{élève 2}& 0{,}3&-0{,}8\\
      \text{élève 3}& 1{,}1&0{,}6
    \end{array}$};
  \node[box,right=30mm of raw] (design) {$
    \X=\begin{bmatrix}
      1&-1{,}2& 0{,}4\\
      1& 0{,}3&-0{,}8\\
      1& 1{,}1& 0{,}6
    \end{bmatrix}$};
  \draw[-{Latex},thick] (raw)--node[above,align=center]{ajout de la\\colonne du biais}(design);
  \coordinate (midboxes) at ($(raw.south)!.5!(design.south)$);
  \node[below=4mm of midboxes,align=center]
    {une ligne $=$ une observation ; une colonne $=$ une caractéristique\\
     la première colonne, constante, est associée à $\theta_0$};
\end{tikzpicture}

\vspace{4mm}
\[
  \underbrace{\begin{bmatrix}
    1&-1{,}2&0{,}4\\1&0{,}3&-0{,}8\\1&1{,}1&0{,}6
  \end{bmatrix}}_{\X\;(3\times3)}
  \underbrace{\begin{bmatrix}\theta_0\\\theta_1\\\theta_2\end{bmatrix}}_{\thv\;(3\times1)}
  =
  \underbrace{\begin{bmatrix}
    \theta_0-1{,}2\theta_1+0{,}4\theta_2\\
    \theta_0+0{,}3\theta_1-0{,}8\theta_2\\
    \theta_0+1{,}1\theta_1+0{,}6\theta_2
  \end{bmatrix}}_{\X\thv=(z_1,z_2,z_3)^\top\;(3\times1)}.
\]
\caption{Construction concrète d'une matrice de design. La colonne de $1$
multiplie $\theta_0$ dans chaque score et encode ainsi le terme constant.
Ici $\X\in\R^{3\times3}$, $\thv\in\R^3$ et $\X\thv\in\R^3$.}
\label{fig:design-matrix-example}
\end{figure}
\FloatBarrier

\begin{procedure}{Former la matrice de design}
\textbf{Entrée :} $m$ observations déjà préparées, chacune décrite par les mêmes
$n$ caractéristiques dans un ordre fixé. \textbf{Sortie :} une matrice
$\X\in\R^{m\times(n+1)}$.
\begin{enumerate}
  \item écrire les $n$ valeurs de l'observation $i$ sur la ligne $i$ ;
  \item ajouter à gauche une colonne de $1$, associée au biais $\theta_0$ ;
  \item enregistrer l'ordre des colonnes, par exemple
        $(1,\text{Astronomy},\text{Herbology},\ldots)$ ;
  \item vérifier que chaque ligne contient $n+1$ nombres et que $\X\thv$ produit
        exactement $m$ scores.
\end{enumerate}
\end{procedure}

Avec deux caractéristiques, le rôle de la constante apparaît dans l'équation de
la frontière. Le score $\theta_1x_1+\theta_2x_2$ place la frontière sur une
droite passant par l'origine. Le score $\theta_0+\theta_1x_1+\theta_2x_2$ la
place sur $\theta_0+\theta_1x_1+\theta_2x_2=0$, qui se déplace parallèlement à
elle-même quand $\theta_0$ varie. Le biais est le paramètre qui règle cette
translation.

\begin{figure}[htbp]
\centering
\begin{tikzpicture}
\begin{axis}[courseplot,axis lines=middle,xlabel={$x_1$},ylabel={$x_2$},
  xmin=-3,xmax=3,ymin=-3,ymax=3,legend style={at={(.5,1.03)},anchor=south,
  legend columns=3,draw=none}]
  \addplot[ink,very thick,domain=-3:3]{-.7*x};
  \addlegendentry{$\theta_0=0$}
  \addplot[accent,very thick,domain=-3:3]{1.3-.7*x};
  \addlegendentry{$\theta_0<0$}
  \addplot[warning,very thick,domain=-3:3]{-1.3-.7*x};
  \addlegendentry{$\theta_0>0$}
  \addplot[only marks,mark=*,ink] coordinates {(0,0)};
\end{axis}
\end{tikzpicture}
\caption{Effet du biais pour des poids $\theta_1$ et $\theta_2$ fixés. Modifier
$\theta_0$ translate la frontière parallèlement à elle-même ; sans biais, elle
est contrainte de passer par l'origine.}
\label{fig:bias-geometry}
\end{figure}

\section{Valeurs manquantes et standardisation}

\begin{lecture}
Le calcul exige un nombre dans chaque case et des échelles comparables. Imputer
consiste à définir la valeur utilisée lorsqu'une mesure manque. Standardiser
consiste à exprimer ensuite chaque mesure par rapport à la moyenne et à
l'écart-type calculés sur l'apprentissage.
\end{lecture}

Les opérations de la régression logistique portent sur des nombres. Chaque case
vide demande donc une règle explicite, par exemple l'imputation médiane, et la
valeur d'imputation de chaque variable rejoint l'état conservé du modèle. La
suppression d'une ligne relève du même principe : un choix, énoncé et documenté.

Exemple : les valeurs d'apprentissage d'un cours sont $\{4,7,9,12,100\}$. La
médiane vaut $9$ et remplace une note manquante. La moyenne vaudrait $26{,}4$,
tirée vers le haut par la seule valeur $100$. La médiane offre donc ici une
robustesse aux valeurs extrêmes ; en contrepartie, l'imputation contracte la
dispersion et absorbe un éventuel mécanisme informatif derrière l'absence.

L'ordre retenu dans ce document --- imputer, puis estimer $\mu_j$ et $s_j$ sur la
colonne complétée --- applique une règle unique à l'apprentissage, à la validation
et à toute observation future. Il contracte $s_j$ d'autant que la proportion
d'imputations est grande. Estimer $\mu_j$ et $s_j$ sur les seules valeurs
observées est l'autre convention défensible. Le point déterminant est
l'enregistrement du choix et son application identique à chaque étape.

\begin{figure}[htbp]
\centering
\begin{tabular}{c@{\hspace{14mm}$\xrightarrow{\text{médiane }=9}$\hspace{14mm}}c}
\begin{tabular}{c}\toprule valeur observée\\\midrule
$4$\\$7$\\\cellcolor{warning!25}\textcolor{warning}{manquante}\\$12$\\$100$\\\bottomrule
\end{tabular}
&
\begin{tabular}{c}\toprule valeur préparée\\\midrule
$4$\\$7$\\\cellcolor{warning!25}$9$\\$12$\\$100$\\\bottomrule
\end{tabular}
\end{tabular}
\caption{Imputation médiane sur une variable. La case colorée porte une valeur
calculée, d'un statut distinct des quatre mesures qui l'entourent. La médiane
provient des seules lignes d'apprentissage.}
\label{fig:median-imputation}
\end{figure}

Pour une note brute $r_j$, on estime sur l'apprentissage
\[
  \mu_j=\frac1{m_{\mathrm{tr}}}\sum_{i\in\mathcal I_{\mathrm{tr}}}r_{ij},
  \qquad
  s_j=\sqrt{\frac1{m_{\mathrm{tr}}}\sum_{i\in\mathcal I_{\mathrm{tr}}}
    (r_{ij}-\mu_j)^2},
  \qquad x_j=\frac{r_j-\mu_j}{s_j}.
\]
Le diviseur $m_{\mathrm{tr}}$ ou $m_{\mathrm{tr}}-1$ change le résultat de moins
d'un pour mille sur $1\,600$ lignes ; le même reste employé à l'entraînement et à
la prédiction. Un $s_j$ nul signale une colonne constante sur l'apprentissage,
que le programme retire de la liste des variables.

La standardisation reparamètre le modèle et conserve la famille des frontières
linéaires. Elle améliore le conditionnement numérique de l'optimisation et rend
une éventuelle pénalisation comparable entre variables
\cite[sec.~4.6.3]{hastie2009}.

Le \emph{conditionnement} mesure la sensibilité du calcul aux directions de
l'espace des paramètres. Une variable qui varie autour de $10^3$ et une autre
autour de $10^{-2}$ répondent à un même déplacement de leurs coefficients par des
variations de score dans un rapport de $10^5$. Les courbes de niveau du coût
s'allongent d'autant, et la descente zigzague ou progresse par pas très courts.
La figure~\ref{fig:contours} en donne la représentation géométrique. La
standardisation ramène les échelles à $1$ et conserve l'ordre des observations
dans chaque colonne.

\begin{resultat}[Exemple de standardisation]
Supposons que la note Astronomy ait, sur l'apprentissage, une moyenne
$\mu=20$ et un écart-type $s=100$. Les notes brutes $-180$, $20$ et $170$
deviennent respectivement $-2$, $0$ et $1{,}5$. La transformation conserve leur
ordre et leurs positions relatives, et fixe une origine et une échelle communes.
Une valeur standardisée de $1{,}5$ se lit « un écart-type et demi au-dessus de la
moyenne d'apprentissage », dans l'unité propre à cette colonne.
\end{resultat}

\begin{figure}[htbp]
\centering
\begin{tikzpicture}[x=1cm,y=1cm]
  \node[anchor=east] at (.4,1.25){note brute $r$};
  \node[anchor=east] at (.4,-.35){valeur $x=(r-20)/100$};
  \draw[{Latex}-{Latex},thick] (.7,1.25)--(11.3,1.25);
  \draw[{Latex}-{Latex},thick] (.7,-.35)--(11.3,-.35);
  \foreach \xpos/\raw/\std in {1/-240/-2{,}6,11/240/2{,}2}{
    \draw (\xpos,1.16)--(\xpos,1.34) node[above=2pt]{\raw};
    \draw (\xpos,-.44)--(\xpos,-.26) node[below=2pt]{\std};
  }
  \foreach \xpos/\raw/\std in {2.25/-180/-2,6.4167/20/0,9.5417/170/1{,}5}{
    \fill[accent] (\xpos,1.25) circle (2.5pt) node[above=7pt]{\raw};
    \fill[warning] (\xpos,-.35) circle (2.5pt) node[below=7pt]{\std};
    \draw[-{Latex},dashed,gridgray!80] (\xpos,1.02)--(\xpos,-.12);
  }
  \node[draw=ink!45,fill=white,rounded corners=2pt,inner sep=4pt]
    at (6,2.35){$r=20+100x$ : même position, autre origine et autre unité};
\end{tikzpicture}
\caption{Correspondance complète entre l'axe des notes brutes et l'axe
standardisé. Chaque flèche verticale relie deux écritures de la même observation :
$-180\leftrightarrow-2$, $20\leftrightarrow0$ et $170\leftrightarrow1{,}5$.}
\label{fig:standardization}
\end{figure}

\begin{vigilance}[Fuite de données]
Calculer $\mu_j$, $s_j$, une médiane d'imputation ou choisir des variables à
partir de l'ensemble complet transmet de l'information du test vers
l'apprentissage, et rend l'évaluation optimiste. La partition précède donc toute
transformation apprise des données.
\end{vigilance}

\begin{procedure}{Préparer une caractéristique numérique}
\textbf{Entrée :} les valeurs brutes $r_{ij}$ de la caractéristique $j$ sur les
lignes d'apprentissage. \textbf{Sortie :} la valeur d'imputation, $\mu_j$, $s_j$
et les valeurs transformées $x_{ij}$.
\begin{enumerate}
  \item calculer la médiane des valeurs présentes et l'utiliser pour remplacer
        les valeurs manquantes ;
  \item calculer ensuite $\mu_j$ et $s_j$ sur cette colonne d'apprentissage complétée ;
  \item transformer chaque valeur par $x_{ij}=(r_{ij}-\mu_j)/s_j$ ; par exemple,
        $r=170$, $\mu=20$, $s=100$ donnent $x=1{,}5$ ;
  \item conserver médiane, $\mu_j$ et $s_j$, et les réappliquer telles quelles à
        la validation, au test et à toute observation future ;
  \item si $s_j=0$, signaler une colonne constante et la retirer du calcul.
\end{enumerate}
\end{procedure}

\section{Partition et déséquilibre des classes}

\begin{lecture}
La partition produit deux listes d'indices disjointes : les indices
d'apprentissage $\mathcal I_{\mathrm{tr}}$ et les indices de test
$\mathcal I_{\mathrm{te}}$. Une ligne appartient à une seule liste. Toutes les
transformations et tous les coefficients sont estimés depuis la première ; la
seconde reste intacte jusqu'à l'évaluation finale.
\end{lecture}

Les deux listes vérifient
\[
  \mathcal I_{\mathrm{tr}}\cap\mathcal I_{\mathrm{te}}=\varnothing,
  \qquad
  \mathcal I_{\mathrm{tr}}\cup\mathcal I_{\mathrm{te}}=\{1,\ldots,m\}.
\]
Pour construire une partition stratifiée, on regroupe d'abord les indices par
classe. Dans chaque groupe, on mélange les indices avec une graine fixée, puis on
affecte la proportion prévue à l'apprentissage et le reste au test. Les morceaux
obtenus sont enfin réunis dans les deux listes. L'opération porte sur les indices
et laisse les données numériques en place.

L'état à conserver tient donc en trois objets : la graine et les deux listes
d'indices. Moyennes, écarts-types et valeurs d'imputation viennent ensuite des
seules lignes désignées par $\mathcal I_{\mathrm{tr}}$, puis s'appliquent telles
quelles aux lignes de $\mathcal I_{\mathrm{te}}$.

\section{Préservation des effectifs par classe}

Avec $1\,000$ élèves, dont $900$ majoritaires et $100$ dans la classe étudiée,
une répartition $80\,\%/20\,\%$ affecte séparément $80\,\%$ de chaque classe à
l'apprentissage. On obtient $720+80=800$ lignes d'apprentissage et
$180+20=200$ lignes de test. Le déséquilibre $90\,\%/10\,\%$ se retrouve à
l'identique dans les deux sous-ensembles.

\begin{figure}[htbp]
\centering
\setlength{\tabcolsep}{12pt}
\renewcommand{\arraystretch}{1.35}
\begin{tabular}{lrrr}
\toprule
classe & ensemble disponible & apprentissage ($80\,\%$) & test ($20\,\%$) \\
\midrule
\textcolor{accent}{majoritaire} & $900$ & $720$ & $180$ \\
\textcolor{warning}{classe étudiée} & $100$ & $80$ & $20$ \\
\midrule
total & $1\,000$ & $800$ & $200$ \\
\bottomrule
\end{tabular}

\vspace{.7em}
\begin{tabular}{ccc}
$900=720+180$ & \qquad & $100=80+20$
\end{tabular}
\caption{Effectifs produits par une partition stratifiée $80\,\%/20\,\%$.}
\label{fig:stratification}
\end{figure}

Un \emph{hyperparamètre} est une quantité que le protocole fixe avant
l'ajustement. Les coefficients $\thv$ sont des paramètres appris par le gradient ;
le taux $\alpha$ et la pénalisation $\lambda$ sont des hyperparamètres. Si trois
valeurs de $\lambda$ donnent sur validation des exactitudes $0{,}78$, $0{,}83$ et
$0{,}80$, la deuxième est retenue, puis le test intervient une seule fois. La
validation croisée du chapitre~\ref{ch:evaluation} étend ce principe en faisant
tourner le rôle de validation sur plusieurs blocs.

\begin{resultat}[Quantités fixées après la préparation]
La préparation produit les listes d'indices d'apprentissage et de test, une
valeur d'imputation par caractéristique, une moyenne et un écart-type par
caractéristique, puis les matrices numériques transformées. Ces quantités sont
estimées une fois sur l'apprentissage et réutilisées sans modification pour la
validation, le test et toute observation future.
\end{resultat}

\begin{procedure}{Construire une partition stratifiée}
\textbf{Entrées :} les numéros de lignes $1,\ldots,m$, leur classe $c_i$, une
proportion d'apprentissage et une \emph{graine}. La graine est un entier qui
initialise le générateur pseudo-aléatoire : reprendre le même entier reproduit
le même mélange. \textbf{Sorties :} deux listes d'indices
$\mathcal I_{\mathrm{tr}}$ et $\mathcal I_{\mathrm{te}}$.
\begin{enumerate}
  \item former une liste de numéros de lignes par classe ; par exemple, une liste
        contient les $100$ indices de la classe étudiée ;
  \item mélanger chaque liste après initialisation avec la graine enregistrée ;
  \item pour une proportion $80\,\%$, placer les $80$ premiers indices de cette
        liste dans $\mathcal I_{\mathrm{tr}}$ et les $20$ autres dans
        $\mathcal I_{\mathrm{te}}$ ; répéter pour chaque classe ;
  \item concaténer les morceaux obtenus : dans l'exemple complet, les deux listes
        contiennent respectivement $800$ et $200$ indices ;
  \item vérifier qu'aucun indice n'apparaît deux fois, qu'aucun ne manque et que
        les effectifs par classe correspondent à la table~\ref{fig:stratification} ;
  \item enregistrer définitivement graine et listes avant de calculer médianes,
        moyennes ou écarts-types.
\end{enumerate}
\end{procedure}

\chapter{Modèle logistique binaire}
\label{ch:binary}

\begin{lecture}
Le cas binaire pose une question à deux réponses, par exemple « Ravenclaw ou
non-Ravenclaw ? ». Les caractéristiques produisent d'abord un score réel, qui
peut être négatif ou positif. La sigmoïde convertit ce score en un nombre entre
$0$ et $1$ interprété comme une probabilité conditionnelle.
\end{lecture}

Considérons provisoirement des étiquettes $y_i\in\{0,1\}$. On pose
\[
  z_{\thv}(\x)=\thv^\top\xt,
  \qquad
  \sigma(z)=\frac1{1+e^{-z}},
  \qquad
  p_{\thv}(\x)=\sigma(z_{\thv}(\x)).
\]
Cette paramétrisation définit la régression logistique binaire
\cite[sec.~4.4]{hastie2009,cox1958}.

Pour un ensemble de $m$ observations, le même calcul est effectué simultanément :
$\X\thv$ produit un vecteur de $m$ scores, puis la sigmoïde appliquée composante
par composante produit un vecteur de $m$ probabilités. Le vecteur $\thv$ contient
les seules quantités ajustées par l'apprentissage ; $\X$ reste fixé pendant les
itérations.

\section{Interprétation par les log-cotes}

\begin{lecture}
Les cotes comparent la probabilité de l'événement à celle de son contraire. Une
probabilité de $0{,}75$ correspond à des cotes $0{,}75/0{,}25=3$ : l'événement
est trois fois plus probable que son contraire. Le logarithme transforme les
multiplications de cotes en additions, compatibles avec un score linéaire.
\end{lecture}

La fonction logistique est équivalente au modèle linéaire
\[
  \log\frac{\Pp(Y=1\mid X=\x)}{\Pp(Y=0\mid X=\x)}
  =\thv^\top\xt.
\]
Cette égalité est précisément la formulation en log-cotes du modèle logistique
\cite[équations~4.17--4.18]{hastie2009}.

Le modèle suppose que les log-cotes sont une fonction affine des
caractéristiques. La transformation inverse est
\[
  p=\sigma(z)=\frac{1}{1+e^{-z}}.
\]
Le calcul reçoit donc un score réel $z=\thv^\top\xt$ et produit une probabilité
comprise entre $0$ et $1$. Les étapes algébriques de l'inversion sont données à
l'annexe~\ref{app:logit-inversion}.
Ainsi, à caractéristiques égales par ailleurs, augmenter $x_j$ d'une unité
multiplie les \emph{cotes} de la classe positive par $e^{\theta_j}$. Après
standardisation, cette unité correspond à un écart-type de la variable brute.
Cette interprétation porte sur les cotes conditionnelles, non sur une variation
constante de probabilité.

Par exemple, $\theta_j=\log 2\approx0{,}693$ signifie qu'une augmentation d'une
unité de $x_j$ double les cotes, puisque $e^{\theta_j}=2$. Si la probabilité
initiale vaut $0{,}20$, les cotes valent $0{,}20/0{,}80=0{,}25$ ; après
doublement, elles valent $0{,}50$, soit une probabilité
$0{,}50/(1+0{,}50)=1/3$. La probabilité n'a donc pas doublé : c'est bien le
rapport de chances qui a été multiplié par deux.

\begin{figure}[htbp]
\centering
\begin{tikzpicture}[node distance=13mm,
  box/.style={draw=accent,fill=accentlight,rounded corners,minimum width=34mm,
    minimum height=13mm,align=center},
  arr/.style={-{Latex},very thick,draw=ink}]
  \node[box] (p1) {probabilité\\$p=0{,}20$};
  \node[box,right=of p1] (o1) {cotes\\$p/(1-p)=0{,}25$};
  \node[box,right=of o1] (l1) {log-cotes\\$\log(0{,}25)=-1{,}386$};
  \draw[arr] (p1)--node[above,scale=.8]{$p/(1-p)$}(o1);
  \draw[arr] (o1)--node[above,scale=.8]{$\log$}(l1);
  \node[box,below=14mm of p1] (p2) {probabilité\\$p=1/3$};
  \node[box,right=of p2] (o2) {cotes doublées\\$0{,}50$};
  \node[box,right=of o2] (l2) {log-cotes augmentées\\de $\log 2$};
  \draw[arr] (p2)--(o2); \draw[arr] (o2)--(l2);
  \draw[warning,very thick,-{Latex}] (o1)--node[left]{${}\times2$}(o2);
  \draw[warning,very thick,-{Latex}] (l1)--node[right]{${}+\log2$}(l2);
\end{tikzpicture}
\caption{Passage entre probabilité, cotes et log-cotes. Ajouter $\log2$ au score
double les cotes, mais transforme ici la probabilité de $0{,}20$ en $1/3$, pas
en $0{,}40$.}
\label{fig:odds-map}
\end{figure}

\begin{figure}[htbp]
  \centering
  \begin{tikzpicture}
    \begin{axis}[courseplot,axis lines=left,
      xlabel={score $z=\thv^\top\xt$},ylabel={$\sigma(z)$},
      xmin=-8,xmax=8,ymin=-.03,ymax=1.03,xtick={-6,-3,0,3,6},
      ytick={0,.25,.5,.75,1},grid=major,grid style={gridgray!50},samples=250]
      \addplot[accent,very thick,domain=-8:8]{1/(1+exp(-x))};
      \addplot[dashed,ink,domain=-8:0]{.5};
      \addplot[dashed,ink] coordinates {(0,0) (0,.5)};
    \end{axis}
  \end{tikzpicture}
  \caption{La logistique est strictement croissante, vérifie
  $\sigma(-z)=1-\sigma(z)$ et $\sigma'(z)=\sigma(z)(1-\sigma(z))$. Le seuil de
  probabilité $1/2$ correspond donc au score nul.}
  \label{fig:sigmoid}
\end{figure}

\begin{figure}[htbp]
\centering
\begin{tikzpicture}
\begin{axis}[courseplot,axis lines=left,
  xlabel={score $z$},ylabel={sortie},xmin=-5,xmax=5,ymin=-.05,ymax=1.05,
  grid=major,grid style={gridgray!45},legend pos=south east,samples=200]
  \addplot[accent,very thick,domain=-5:5]{1/(1+exp(-x))};
  \addlegendentry{probabilité $\sigma(z)$}
  \addplot[warning,very thick,const plot] coordinates {(-5,0) (0,0) (0,1) (5,1)};
  \addlegendentry{décision $\ind[z\geq0]$}
\end{axis}
\end{tikzpicture}
\caption{La probabilité et la décision ne sont pas le même objet. La sigmoïde
exprime une transition continue et conserve le degré de confiance ; le seuil
produit seulement une étiquette binaire discontinue.}
\label{fig:sigmoid-threshold}
\end{figure}

\begin{procedure}{Calculer une sortie binaire}
\textbf{Entrées :} une ligne $\xt\in\R^{n+1}$ et les paramètres
$\thv\in\R^{n+1}$. \textbf{Sorties :} un score réel $z$ et une probabilité $p$.
\begin{enumerate}
  \item vérifier que $\xt$ et $\thv$ ont la même longueur et le même ordre de colonnes ;
  \item calculer $z=\theta_0+\theta_1x_1+\cdots+\theta_nx_n$ ;
  \item calculer $p=1/(1+e^{-z})$ avec la forme stable du chapitre~\ref{ch:optim} ;
  \item vérifier $0<p<1$ et conserver $z$ séparément de $p$ : OvR compare les scores.
\end{enumerate}
\end{procedure}

\section{Frontière de décision}

\begin{lecture}
La probabilité varie continûment, mais la décision exige une classe. Avec un
seuil de $0{,}5$, la frontière réunit les observations pour lesquelles le modèle
hésite exactement entre les deux réponses. Dans deux dimensions, cette frontière
est une droite ; dans davantage de dimensions, c'est un hyperplan.
\end{lecture}

Avec des coûts d'erreur symétriques, la règle de Bayes estimée est
\[
  \widehat y(\x)=\ind[p_{\thv}(\x)\geq 1/2]
  =\ind[\thv^\top\xt\geq0].
\]
La règle générale consiste à choisir la classe de probabilité conditionnelle
maximale ; le seuil $1/2$ en est le cas binaire à coûts symétriques
\cite[sec.~2.4]{hastie2009}.
La frontière $\{\x:\thv^\top\xt=0\}$ est un hyperplan. La transformation
logistique rend la sortie bornée, mais n'introduit aucune non-linéarité dans la
géométrie de la frontière exprimée dans les variables fournies au modèle.

\begin{figure}[htbp]
\centering
\begin{tikzpicture}
\begin{axis}[courseplot,axis lines=middle,
  xlabel={$x_1$},ylabel={$x_2$},xmin=-3,xmax=3,ymin=-3,ymax=3,
  grid=major,grid style={gridgray!45},legend pos=south east]
  \addplot[only marks,mark=*,accent,mark size=2.4pt]
    coordinates {(-2,-1.2) (-1.5,-.4) (-1.8,.5) (-.8,-1.4) (-.4,-.8)};
  \addlegendentry{$y=0$}
  \addplot[only marks,mark=triangle*,warning,mark size=2.8pt]
    coordinates {(.3,1.2) (1,.6) (1.4,1.8) (2,.8) (1.6,-.2)};
  \addlegendentry{$y=1$}
  \addplot[ink,very thick,domain=-3:3]{.4-.75*x};
  \addlegendentry{$\theta_0+\theta_1x_1+\theta_2x_2=0$}
  \addplot[accent!55,dashed,domain=-3:3]{1.5-.75*x};
  \addplot[accent!55,dashed,domain=-3:3]{-0.7-.75*x};
\end{axis}
\end{tikzpicture}
\caption{Exemple à deux caractéristiques. La droite pleine est la frontière
$p=0{,}5$ ; les droites pointillées sont des lignes d'égale probabilité. Elles
sont parallèles parce que le score est affine. Les points sont illustratifs et
ne représentent pas les données complètes du projet.}
\label{fig:boundary}
\end{figure}

\begin{procedure}{Appliquer une décision binaire}
\textbf{Entrée :} une probabilité $p$ et un seuil $\tau$ fixé avant le test.
\textbf{Sortie :} une étiquette $\widehat y\in\{0,1\}$.
\begin{enumerate}
  \item avec des coûts symétriques, choisir $\tau=0{,}5$ ;
  \item produire $\widehat y=1$ si $p\geq\tau$, sinon $\widehat y=0$ ;
  \item exemple : $p=0{,}71$ donne $1$, tandis que $p=0{,}23$ donne $0$ ;
  \item si un autre seuil est nécessaire, le sélectionner sur validation et
        enregistrer sa valeur avant l'évaluation du test.
\end{enumerate}
\end{procedure}

\begin{resultat}[Quantités produites par un modèle binaire]
Une ligne préparée $\xt$ et un vecteur de paramètres $\thv$ produisent d'abord
un score $z=\thv^\top\xt$, puis une probabilité $p=\sigma(z)$, puis une décision
binaire après comparaison au seuil. L'apprentissage modifie $\thv$ ; la
prédiction utilise ensuite ce vecteur sans le réestimer.
\end{resultat}

\chapter{Estimation par maximum de vraisemblance}
\label{ch:likelihood}

\begin{lecture}
L'apprentissage compare les probabilités annoncées aux réponses observées. Un
paramétrage reçoit une vraisemblance élevée lorsqu'il attribue une forte
probabilité aux étiquettes effectivement présentes dans le fichier. La
log-perte est la même comparaison écrite sous la forme d'un coût à minimiser.
\end{lecture}

On suppose les observations indépendantes conditionnellement à leurs variables
explicatives et
\[
  Y_i\mid X_i=\x_i\sim
  \operatorname{Bernoulli}(p_i),\qquad p_i=\sigma(\thv^\top\xt_i).
\]
Ce modèle de Bernoulli conditionnel et son estimation par maximum de
vraisemblance sont la construction standard de la régression logistique
\cite[sec.~4.4.1]{hastie2009,cox1958}.

À une valeur donnée de $\thv$ correspondent donc trois objets successifs : le
vecteur des scores $\X\thv$, le vecteur des probabilités $\mathbf p$ et un coût scalaire
$J(\thv)$ qui résume l'accord avec toutes les cibles. L'apprentissage réévalue
ces objets après chaque modification de $\thv$ ; les observations et leurs
cibles restent fixes.

La masse de probabilité d'une Bernoulli peut s'écrire sous une forme unique pour
les deux valeurs possibles de $y_i$ :
\[
  \Pp(Y_i=y_i\mid X_i=\x_i;\thv)
  =p_i^{y_i}(1-p_i)^{1-y_i}.
\]
Si $y_i=1$, cette expression devient $p_i$ ; si $y_i=0$, elle devient $1-p_i$.
Sous l'hypothèse d'indépendance conditionnelle des observations, la probabilité
jointe des réponses observées est le produit de ces contributions. C'est ce
passage --- et non une analogie avec les moindres carrés --- qui produit la
vraisemblance suivante.

Cette indépendance concerne les élèves une fois leurs caractéristiques fixées.
Elle permet d'écrire, pour deux observations,
\begin{align*}
  &\Pp(Y_1=y_1,Y_2=y_2\mid X_1=\x_1,X_2=\x_2;\thv)\\
  &\qquad=\Pp(Y_1=y_1\mid X_1=\x_1;\thv)
   \Pp(Y_2=y_2\mid X_2=\x_2;\thv).
\end{align*}
Elle serait douteuse si plusieurs lignes correspondaient au même élève ou à des
jumeaux dont les réponses restent liées après prise en compte des variables. Le
produit de vraisemblance n'est donc pas une identité universelle : il traduit une
hypothèse sur la manière dont les données ont été recueillies.
La vraisemblance et sa log-vraisemblance sont
\begin{align*}
  L(\thv)&=\prod_{i=1}^m p_i^{y_i}(1-p_i)^{1-y_i},\\
  \log L(\thv)&=\sum_{i=1}^m
  \left[y_i\log p_i+(1-y_i)\log(1-p_i)\right].
\end{align*}
Maximiser $L$ équivaut à minimiser la log-perte moyenne
\[
  J(\thv)=-\frac1m\log L(\thv)
  =-\frac1m\sum_{i=1}^m\left[y_i\log p_i+
  (1-y_i)\log(1-p_i)\right].
  \tag{1}\label{eq:loss}
\]
L'équation~\eqref{eq:loss} est l'opposé de la log-vraisemblance binomiale,
également appelée log-perte ou entropie croisée binaire
\cite[équations~4.19--4.20]{hastie2009}.
Cette quantité est l'opposé du logarithme de la probabilité attribuée par le
modèle aux étiquettes observées.

\begin{table}[htbp]
\centering
\begin{tabular}{ccccc}
\toprule
$i$ & $y_i$ & $p_i$ & probabilité attribuée au résultat observé & $\ell_i$ \\
\midrule
$1$ & $1$ & $0{,}90$ & $0{,}90$ & $-\log(0{,}90)=0{,}105$ \\
$2$ & $0$ & $0{,}20$ & $1-0{,}20=0{,}80$ & $-\log(0{,}80)=0{,}223$ \\
$3$ & $1$ & $0{,}10$ & $0{,}10$ & $-\log(0{,}10)=2{,}303$ \\
\midrule
\multicolumn{4}{r}{perte moyenne $J$} & $0{,}877$ \\
\bottomrule
\end{tabular}
\caption{Calcul de la log-perte sur trois observations. La troisième prédiction,
fausse et assurée, domine la moyenne ; c'est le comportement recherché par le
maximum de vraisemblance.}
\label{tab:loss-example}
\end{table}

\begin{figure}[htbp]
  \centering
  \begin{tikzpicture}
    \begin{axis}[courseplot,axis lines=left,
      xlabel={probabilité prédite $p$},ylabel={perte $\ell(y,p)$},
      xmin=0,xmax=1,ymin=0,ymax=5,grid=major,grid style={gridgray!50},
      samples=250,domain=.007:.993,legend pos=north west]
      \addplot[accent,very thick]{-ln(x)}; \addlegendentry{$y=1$: $-\log p$}
      \addplot[warning,very thick]{-ln(1-x)}; \addlegendentry{$y=0$: $-\log(1-p)$}
    \end{axis}
  \end{tikzpicture}
  \caption{Une prédiction correcte et assurée coûte presque zéro ; une
  prédiction erronée et assurée reçoit une pénalité non bornée.}
  \label{fig:loss}
\end{figure}

\begin{figure}[htbp]
\centering
\begin{tikzpicture}
\begin{axis}[courseplot,axis lines=left,
  xlabel={probabilité prédite $p$ pour une observation $y=1$},ylabel={pénalité},
  xmin=.02,xmax=1,ymin=0,ymax=4.2,grid=major,grid style={gridgray!45},
  legend pos=north east,samples=240,domain=.02:1]
  \addplot[warning,very thick]{-ln(x)};
  \addlegendentry{log-perte $-\log p$}
  \addplot[accent,dashed,very thick]{(1-x)^2};
  \addlegendentry{erreur quadratique $(1-p)^2$}
\end{axis}
\end{tikzpicture}
\caption{Comparaison pour $y=1$. Les deux pertes sont minimales en $p=1$, mais
la log-perte pénalise beaucoup plus fortement une probabilité presque nulle
attribuée à l'événement réellement observé. Elle découle en outre directement
de la vraisemblance de Bernoulli.}
\label{fig:loss-comparison}
\end{figure}

\begin{procedure}{Évaluer la log-perte}
\textbf{Entrées :} $m$ probabilités $p_i$ et $m$ cibles $y_i\in\{0,1\}$.
\textbf{Sortie :} un nombre $J(\thv)$ ; une valeur plus faible indique un meilleur
accord avec les étiquettes observées.
\begin{enumerate}
  \item calculer les scores $z_i$, puis les probabilités $p_i$ ;
  \item pour chaque ligne, prendre $a_i=p_i$ si $y_i=1$, sinon $a_i=1-p_i$ ;
  \item calculer $\ell_i=-\log(a_i)$ ; par exemple, $a_i=0{,}9$ donne
        $\ell_i\approx0{,}105$ et $a_i=0{,}1$ donne $\ell_i\approx2{,}303$ ;
  \item calculer $J=m^{-1}\sum_i\ell_i$ et vérifier qu'il reste fini.
\end{enumerate}
\end{procedure}

\section{Gradient et correction des coefficients}

\begin{lecture}
Le gradient indique comment une petite modification de chaque coefficient ferait
varier le coût. Chaque observation contribue selon son erreur $p_i-y_i$ et selon
ses caractéristiques. La somme de ces contributions fournit la correction de
tous les coefficients en une opération matricielle.
\end{lecture}

Pour chaque observation, le programme calcule d'abord l'écart $p_i-y_i$ entre
la probabilité produite et la cible binaire. Il pondère ensuite cet écart par les
caractéristiques de l'observation, puis additionne les contributions. Le calcul
complet s'écrit
\[
  \boxed{\quad\nabla J(\thv)=\frac1m\X^\top
  \bigl(\sigma(\X\thv)-\y\bigr).\quad}
  \tag{2}\label{eq:gradient}
\]
Cette expression est le score de vraisemblance sous forme matricielle
\cite[équation~4.21]{hastie2009}.
Sa dérivation par la règle de chaîne figure à
l'annexe~\ref{app:gradient-derivation}.
Le gradient a bien la dimension $n+1$ : $\X^\top$ est de taille
$(n+1)\times m$ et le vecteur des résidus probabilistes est de taille $m$.

Le terme $(p_i-y_i)\xt_i$ permet une lecture locale. Pour un positif
$y_i=1$ sous-estimé à $p_i=0{,}1$, le résidu vaut $-0{,}9$ : l'observation exerce
une correction importante. Pour un positif déjà estimé à $0{,}99$, le résidu
vaut seulement $-0{,}01$. Le gradient agrège ces corrections en tenant compte de
la direction $\xt_i$ de chaque observation.

\begin{procedure}{Calculer le gradient}
\textbf{Entrées :} $\X\in\R^{m\times(n+1)}$, $\y\in\{0,1\}^m$ et
$\thv\in\R^{n+1}$. \textbf{Sortie :} $\mathbf g\in\R^{n+1}$, dont chaque
composante indique la correction associée à un coefficient.
\begin{enumerate}
  \item calculer $\mathbf z=\X\thv$, de longueur $m$, puis
        $\mathbf p=\sigma(\mathbf z)$ ;
  \item calculer $\mathbf r=\mathbf p-\y$ ; un positif prédit à $0{,}1$ contribue
        par un écart $-0{,}9$ ;
  \item calculer $\mathbf g=\X^\top\mathbf r/m$ ;
  \item vérifier que $\mathbf g$ et $\thv$ ont tous deux $n+1$ composantes ;
  \item avant l'entraînement complet, comparer quelques composantes à la
        différence finie de la figure~\ref{fig:finite-difference}.
\end{enumerate}
\end{procedure}

\section{Convexité, unicité et séparabilité}

\begin{lecture}
Une fonction convexe a la géométrie générale d'une cuvette : descendre localement
ne conduit pas vers un faux minimum isolé. Une cuvette peut toutefois posséder un
fond plat, donc plusieurs paramètres équivalents. Avec des classes parfaitement
séparées, le minimum non pénalisé peut même être repoussé à l'infini.
\end{lecture}

En notant $W=\operatorname{diag}(p_i(1-p_i))$, la Hessienne vaut
\[
  \nabla^2J(\thv)=\frac1m\X^\top W\X.
\]
Cette Hessienne est la matrice d'information observée, à un facteur de
normalisation près \cite[équations~4.22--4.25]{hastie2009}.
Les coefficients diagonaux $p_i(1-p_i)$ sont positifs ou nuls ; la Hessienne est
donc positive semi-définie et $J$ est convexe. Tout minimum local est global. La
preuve matricielle est donnée à l'annexe~\ref{app:convexity-proof}. La stricte
convexité, et donc l'unicité du minimiseur, exige toutefois des conditions de rang
et l'absence de directions dégénérées.

La distinction entre convexité et stricte convexité est importante. Une fonction
convexe ne possède pas de minimum local trompeur, mais elle peut présenter un
fond plat contenant plusieurs minimiseurs. Si deux colonnes de $\X$ sont
identiques, par exemple, certains changements opposés de leurs coefficients
laissent les scores inchangés. La prédiction peut alors être unique tandis que
le vecteur de paramètres ne l'est pas.

En une dimension, $J(\theta)=(\theta-1)^2$ est strictement convexe : sa courbure
$J''(\theta)=2$ est positive et son unique minimum est $\theta=1$. La fonction
$J(\theta_1,\theta_2)=(\theta_1+\theta_2-1)^2$ est seulement convexe : tous les
couples vérifiant $\theta_1+\theta_2=1$ minimisent le coût. Cette seconde
situation reproduit ce qui se passe lorsque deux colonnes de la matrice de design
apportent exactement la même information.

\begin{vigilance}[Séparation complète]
Si un hyperplan sépare parfaitement les deux classes, la vraisemblance non
pénalisée peut croître lorsque $\|\thv\|\to\infty$ sans atteindre de maximiseur
fini. La convexité seule n'assure donc pas l'existence d'un minimiseur fini. Une
pénalisation $L^2$, ou une autre méthode adaptée, rétablit en
général un problème bien posé \cite{albert1984}.
\end{vigilance}

Exemple : sur une caractéristique $x$, supposons que tous les négatifs vérifient
$x<0$ et tous les positifs $x>0$. Le score $z=10x$ classe déjà les observations
avec assurance ; $z=100x$ rend les probabilités encore plus proches de $0$ ou de
$1$ et augmente encore la vraisemblance. Aucun coefficient fini n'est désigné
comme optimum : le modèle pousse le coefficient vers $+\infty$. La pénalisation
ajoute un coût croissant avec ce coefficient et empêche cette fuite.

\begin{figure}[htbp]
\centering
\begin{tikzpicture}
\begin{axis}[courseplot,height=4.8cm,axis x line=middle,axis y line=none,
  xlabel={$x$},xmin=-3,xmax=3,ymin=-.25,ymax=.45,ytick=\empty,
  xtick={-3,-2,-1,0,1,2,3},
  legend style={at={(.5,1.02)},anchor=south,legend columns=2,draw=none}]
  \addplot[only marks,mark=*,accent,mark size=2.8pt]
    coordinates {(-2.5,0) (-1.8,0) (-1.2,0) (-.5,0)};
  \addlegendentry{$y=0$}
  \addplot[only marks,mark=triangle*,warning,mark size=3.2pt]
    coordinates {(.4,0) (1.1,0) (1.7,0) (2.5,0)};
  \addlegendentry{$y=1$}
  \draw[ink,very thick] (axis cs:0,-.09)--(axis cs:0,.13)
    node[above=2pt]{seuil $x=0$};
\end{axis}
\end{tikzpicture}
\caption{Séparation complète sur une unique caractéristique $x$. Tous les points
appartiennent au même axe : les cercles ($y=0$) sont à gauche du seuil et les
triangles ($y=1$) à droite. Multiplier le coefficient conserve le seuil mais rend
les probabilités toujours plus extrêmes.}
\label{fig:complete-separation}
\end{figure}

\begin{procedure}{Diagnostiquer l'existence d'une solution finie}
\textbf{Entrées :} la matrice $\X$, les cibles $\y$ et l'historique des
paramètres. \textbf{Sortie :} un diagnostic : ajustement ordinaire, redondance
des colonnes ou séparation complète.
\begin{enumerate}
  \item repérer deux colonnes identiques ou proportionnelles : leurs coefficients
        ne peuvent pas être identifiés séparément ;
  \item vérifier si une frontière classe toutes les observations sans erreur ;
  \item pendant l'ajustement, suivre $\|\thv\|_2$ en plus de $J$ : une norme qui
        croît continuellement tandis que $J$ baisse signale une séparation probable ;
  \item dans ce cas, activer une pénalisation $L^2$ et documenter le diagnostic.
\end{enumerate}
\end{procedure}

\section{Régularisation facultative}

\begin{lecture}
La régularisation compare deux objectifs : bien reproduire les observations et
éviter des coefficients excessivement grands. Le paramètre $\lambda$ fixe le
poids du second objectif. Une valeur nulle conserve le modèle initial ; une
valeur croissante contracte progressivement les coefficients vers zéro.
\end{lecture}

Une pénalisation quadratique donne
\[
  J_\lambda(\thv)=J(\thv)+\frac{\lambda}{2m}\sum_{j=1}^n\theta_j^2,
  \qquad
  \nabla J_\lambda=\nabla J+\frac{\lambda}{m}(0,\theta_1,\ldots,\theta_n)^\top.
\]
Cette forme correspond à la régression logistique pénalisée par une norme
quadratique \cite[sec.~4.4.4]{hastie2009}. Le symbole
\[
  \|\mathbf w\|_2^2=\theta_1^2+\cdots+\theta_n^2
\]
est le carré de la norme euclidienne des coefficients hors biais. Ajouter cette
quantité au coût signifie qu'un ajustement n'est plus jugé uniquement sur son
accord aux données : entre deux solutions de log-pertes voisines, on préfère
celle dont les coefficients sont moins grands.

Le paramètre $\lambda$ règle le compromis. Lorsque $\lambda=0$, on retrouve le
maximum de vraisemblance non pénalisé. Quand $\lambda$ augmente, les coefficients
sont progressivement ramenés vers zéro. À la limite, le modèle dépend de moins en
moins des caractéristiques et tend vers une prédiction principalement gouvernée
par le biais. Le biais n'est conventionnellement pas pénalisé : même si toutes
les caractéristiques deviennent sans effet, le modèle doit encore pouvoir
représenter la fréquence moyenne de la classe positive.

La réduction de variance signifie ici que de petites
modifications de l'échantillon d'apprentissage provoquent généralement de moins
grandes modifications des coefficients et des prédictions. Cette stabilité est
achetée au prix d'un biais statistique supplémentaire : le modèle pénalisé
n'essaie volontairement pas d'ajuster aussi étroitement les données observées.
Le meilleur compromis n'est pas connu à l'avance.

\begin{figure}[htbp]
\centering
\begin{tikzpicture}
\begin{axis}[courseplot,xmode=log,
  xlabel={intensité de régularisation $\lambda$ (échelle logarithmique)},
  ylabel={valeur du coefficient},xmin=.01,xmax=100,ymin=-2.2,ymax=2.8,
  grid=major,grid style={gridgray!40},legend style={at={(.5,1.03)},
  anchor=south,legend columns=3,draw=none},samples=160]
  \addplot[accent,very thick,domain=.01:100]{2.5/(1+x^.55)};
  \addlegendentry{$\theta_1$}
  \addplot[warning,very thick,domain=.01:100]{-1.9/(1+x^.48)};
  \addlegendentry{$\theta_2$}
  \addplot[ink,dashed,very thick,domain=.01:100]{1.25/(1+x^.7)};
  \addlegendentry{$\theta_3$}
  \addplot[gridgray,thin,domain=.01:100]{0};
\end{axis}
\end{tikzpicture}
\caption{Chemins de coefficients sous pénalisation $L^2$ : leur amplitude
décroît lorsque $\lambda$ augmente.}
\label{fig:l2-path}
\end{figure}
\FloatBarrier

On compare donc plusieurs valeurs de $\lambda$ sur un ensemble de validation ou
par validation croisée. Choisir $\lambda$ sur le test reviendrait à adapter le
modèle aux réponses du test ; celui-ci ne fournirait alors plus une évaluation
indépendante. La figure~\ref{fig:lambda-validation} résume cette logique.

\begin{figure}[htbp]
\centering
\begin{tikzpicture}
\begin{axis}[courseplot,xmode=log,
  xlabel={$\lambda$},ylabel={log-perte},xmin=.01,xmax=100,ymin=.2,ymax=1.05,
  grid=major,grid style={gridgray!40},legend pos=north west,samples=180]
  \addplot[accent,very thick,domain=.01:100]{.25+.12*ln(1+x)};
  \addlegendentry{apprentissage}
  \addplot[warning,very thick,domain=.01:100]{.43+.055*(ln(x)-.55)^2};
  \addlegendentry{validation}
  \draw[dashed,ink] (axis cs:1.73,.2)--(axis cs:1.73,.9)
    node[pos=.8,right]{$\lambda$ retenu};
\end{axis}
\end{tikzpicture}
\caption{Sélection de $\lambda$ par minimisation de la log-perte de validation ;
le test final reste exclu de cette sélection.}
\label{fig:lambda-validation}
\end{figure}

\begin{procedure}{Sélectionner la pénalisation}
\textbf{Entrées :} une grille explicite, par exemple
$\lambda\in\{0,10^{-2},10^{-1},1,10\}$, et une partition de validation commune.
\textbf{Sortie :} une valeur de $\lambda$ choisie avant le test.
\begin{enumerate}
  \item réinitialiser et ajuster un modèle pour chaque $\lambda$ sur les mêmes
        lignes d'apprentissage ;
  \item calculer pour chacun la log-perte sur les mêmes lignes de validation ;
  \item construire le tableau $(\lambda,J_{\mathrm{val}})$ et choisir la plus
        petite perte, avec une règle de départage fixée en cas d'égalité ;
  \item enregistrer $\lambda$ et les autres hyperparamètres retenus ;
  \item appliquer ensuite le protocole final et consulter le test une seule fois.
\end{enumerate}
\end{procedure}

\begin{resultat}[Quantités recalculées pendant l'ajustement]
À chaque itération, les paramètres courants produisent les scores, les
probabilités, la log-perte moyenne et le gradient. La log-perte mesure l'état de
l'ajustement ; le gradient produit la correction suivante. La régularisation
ajoute sa contribution au coût et au gradient avant cette correction.
\end{resultat}

\chapter{Optimisation et calcul numérique}
\label{ch:optim}

\begin{lecture}
La fonction de coût associe un nombre à chaque choix des coefficients. Optimiser
consiste à déplacer progressivement ces coefficients vers une zone où ce nombre
est plus faible. Le gradient donne la direction de plus forte hausse ; la
descente suit donc la direction opposée, avec une longueur de pas réglée par
$\alpha$.
\end{lecture}

La descente de gradient simultanée produit la suite
\[
  \thv^{[t+1]}=\thv^{[t]}-\alpha\nabla J(\thv^{[t]}),
  \qquad \alpha>0.
\]
Cette itération est la descente de gradient par lot appliquée à la
log-vraisemblance négative \cite[chap.~8]{boyd2004}.

Le gradient indique la direction locale de hausse du coût ; son opposé indique
donc une direction locale de baisse. Cette propriété résulte du développement
au premier ordre présenté à l'annexe~\ref{app:descent-direction}. Elle reste
locale : un pas trop grand peut dépasser la zone où cette approximation décrit
correctement la fonction.
Le mot \emph{simultanée} est important : toutes les composantes du nouveau
vecteur sont calculées à partir du même ancien vecteur. Un taux trop grand peut
faire augmenter le coût ou provoquer une divergence ; un taux très petit rend
la convergence inutilement lente.

\begin{figure}[htbp]
  \centering
  \begin{tikzpicture}
    \begin{axis}[courseplot,axis lines=left,
      xlabel={itération $t$},ylabel={$J(\thv^{[t]})$},xmin=0,xmax=30,
      ymin=.18,ymax=1.05,grid=major,grid style={gridgray!50},samples=120]
      \addplot[accent,very thick,domain=0:30]{.22+.75*exp(-x/5)};
      \addplot[warning,dashed,thick,domain=0:30]{.32+.55*exp(-x/12)+.08*sin(deg(x))};
    \end{axis}
  \end{tikzpicture}
  \caption{Diagnostic typique du coût : décroissance stable (trait plein) et
  oscillations dues à un pas trop agressif (pointillé). Cette figure est
  schématique ; le tracé du coût réel doit être conservé avec l'expérience.}
\end{figure}

Un arrêt robuste combine un nombre maximal d'itérations avec, par exemple,
\[
  \frac{|J^{(t+1)}-J^{(t)}|}{\max(1,|J^{(t)}|)}<\eps
  \quad\text{ou}\quad \|\nabla J(\thv^{[t]})\|_2<\eps.
\]
Un simple nombre fixe d'itérations ne démontre pas la convergence.

\begin{figure}[htbp]
\centering
\begin{tikzpicture}
\begin{axis}[width=.92\textwidth,height=11.2cm,
  xlabel={$\theta_0$},ylabel={$\theta_1$},xmin=-3,xmax=3,ymin=-3,ymax=3,
  axis equal image,grid=major,grid style={gridgray!30}]
  \foreach \a/\b in {.35/.65,.65/1.2,1/1.85,1.35/2.5}{
    \addplot[accent!55,thick,domain=0:360,samples=140]
      ({-.30+\a*cos(x)},{.35+\b*sin(x)});
  }
  \addplot[warning,very thick,mark=*,mark size=2.1pt,-{Latex}]
    coordinates {(-2.55,2.55) (-1.65,1.92) (-1.05,1.32) (-.68,.86)
      (-.46,.56) (-.34,.40) (-.30,.35)};
  \addplot[only marks,mark=star,mark size=5pt,ink]
    coordinates {(-.30,.35)} node[above right]{minimum};
\end{axis}
\end{tikzpicture}
\caption{Courbes de niveau schématiques de $J(\theta_0,\theta_1)$ et trajectoire
de descente. L'allongement des ellipses traduit un mauvais conditionnement : un
même taux d'apprentissage progresse vite dans une direction et lentement dans
l'autre. Les points successifs aboutissent ici au centre des courbes, le minimum.
La standardisation tend à réduire cet effet d'allongement.}
\label{fig:contours}
\end{figure}

Il faut distinguer trois diagnostics. Une baisse du coût indique que l'algorithme
avance ; une petite norme du gradient indique qu'il approche un point
stationnaire ; une bonne performance de validation indique que ce point produit
un modèle utile hors échantillon. Aucun de ces faits n'implique automatiquement
les deux autres.

Un \emph{point stationnaire} est un point où le gradient est nul. Pour
$J(\theta)=\theta^2$, le point $0$ est stationnaire et constitue le minimum. Pour
$J(\theta)=\theta^3$, le point $0$ est aussi stationnaire, mais la fonction
continue de décroître à gauche et de croître à droite : ce n'est pas un minimum.
Dans le problème logistique convexe, un point stationnaire fini est un minimum
global ; c'est la convexité établie au chapitre précédent qui autorise cette
conclusion, pas la seule égalité $\nabla J=0$.

\paragraph{Influence du taux d'apprentissage dans le cas scalaire.}
La fonction $J(\theta)=(\theta-2)^2$ mesure ici la distance quadratique au
minimum situé en $\theta=2$. À chaque tour, le programme évalue la pente
$J'(\theta)=2(\theta-2)$, la multiplie par $\alpha$, puis soustrait cette
correction à la valeur courante. Depuis $\theta^{[0]}=-2$, le pas
$\alpha=0{,}25$ rapproche successivement le paramètre du minimum, tandis que
$\alpha=1$ le fait passer alternativement d'un côté à l'autre sans réduire
l'écart.

\begin{figure}[htbp]
\centering
\begin{minipage}[t]{.48\textwidth}
\centering
\begin{tikzpicture}
\begin{axis}[width=.85\linewidth,height=6.2cm,axis lines=left,
  title={$\alpha=0{,}25$ : rapprochement du minimum},
  xlabel={itération $t$},ylabel={paramètre $\theta^{[t]}$},
  xmin=0,xmax=6,ymin=-2.6,ymax=6.6,xtick={0,1,2,3,4,5,6},
  grid=major,grid style={gridgray!45}]
  \addplot[accent,very thick,mark=*] coordinates
    {(0,-2) (1,0) (2,1) (3,1.5) (4,1.75) (5,1.875) (6,1.9375)};
  \addplot[ink,dashed,thick,domain=0:6]{2};
  \node[anchor=south east] at (axis cs:6,2){minimum $\theta=2$};
\end{axis}
\end{tikzpicture}
\end{minipage}\hfill
\begin{minipage}[t]{.48\textwidth}
\centering
\begin{tikzpicture}
\begin{axis}[width=.85\linewidth,height=6.2cm,axis lines=left,
  title={$\alpha=1$ : oscillation persistante},
  xlabel={itération $t$},ylabel={paramètre $\theta^{[t]}$},
  xmin=0,xmax=6,ymin=-2.6,ymax=6.6,xtick={0,1,2,3,4,5,6},
  grid=major,grid style={gridgray!45}]
  \addplot[warning,very thick,mark=*] coordinates
    {(0,-2) (1,6) (2,-2) (3,6) (4,-2) (5,6) (6,-2)};
  \addplot[ink,dashed,thick,domain=0:6]{2};
  \node[anchor=south east] at (axis cs:6,2){minimum $\theta=2$};
\end{axis}
\end{tikzpicture}
\end{minipage}
\caption{Valeur du paramètre après chaque mise à jour. Les segments relient deux
itérations successives ; ils représentent l'ordre des calculs, pas une grandeur
supplémentaire.}
\label{fig:learning-rate-scalar}
\end{figure}

Dans le modèle complet, la même séquence s'applique à un vecteur : calculer les
scores, les probabilités, le gradient, puis mettre à jour simultanément tous les
coefficients. Le suivi de $J$ permet de détecter une décroissance, une oscillation
ou une divergence.

\begin{figure}[htbp]
\centering
\resizebox{\textwidth}{!}{%
\begin{tikzpicture}[
  b/.style={draw=accent,fill=accentlight,rounded corners=2pt,
    minimum width=27mm,minimum height=11mm,align=center},
  a/.style={-{Latex},thick,draw=ink},node distance=8mm]
  \node[b] (input) {$\X,\y,\thv^{[t]}$\\données et paramètres};
  \node[b,right=of input] (score) {$\mathbf z=\X\thv^{[t]}$\\scores};
  \node[b,right=of score] (prob) {$\mathbf p=\sigma(\mathbf z)$\\probabilités};
  \node[b,right=of prob] (res) {$\mathbf r=\mathbf p-\y$\\écarts};
  \node[b,right=of res] (grad) {$\mathbf g=\X^\top\mathbf r/m$\\gradient};
  \node[b,below=13mm of grad] (update)
    {$\thv^{[t+1]}=\thv^{[t]}-\alpha\mathbf g$\\mise à jour simultanée};
  \node[b,left=of update] (control) {$J,\|\mathbf g\|_2$\\contrôles d'arrêt};
  \draw[a] (input)--(score); \draw[a] (score)--(prob);
  \draw[a] (prob)--(res); \draw[a] (res)--(grad);
  \draw[a] (grad)--(update); \draw[a] (update)--(control);
  \draw[a] (control.west)-| node[pos=.25,below]{itération suivante} (input.south);
\end{tikzpicture}}
\caption{Cycle de calcul d'une itération de descente de gradient par lot. Chaque
flèche correspond à une transformation nécessaire ; aucun réajustement des
données préparées n'intervient dans cette boucle.}
\label{fig:training-cycle}
\end{figure}

\begin{procedure}{Exécuter une itération d'ajustement}
\textbf{Entrées :} $\X$, $\y$, les paramètres courants $\thv^{[t]}$ et le taux
$\alpha$. \textbf{Sorties :} $\thv^{[t+1]}$, le coût $J^{[t]}$ et la norme du
gradient $\|\mathbf g^{[t]}\|_2$.
\begin{enumerate}
  \item calculer $\mathbf z=\X\thv^{[t]}$, $\mathbf p=\sigma(\mathbf z)$,
        $J^{[t]}$ et $\mathbf g^{[t]}=\X^\top(\mathbf p-\y)/m$ ;
  \item calculer toutes les composantes de
        $\thv^{[t+1]}=\thv^{[t]}-\alpha\mathbf g^{[t]}$ depuis le même ancien vecteur ;
  \item ajouter $J^{[t]}$ et $\|\mathbf g^{[t]}\|_2$ aux historiques ;
  \item arrêter si la tolérance ou le nombre maximal d'itérations est atteint ;
  \item sinon remplacer $t$ par $t+1$ ; si $J$ augmente durablement ou devient
        non fini, interrompre et revoir $\alpha$ ou la stabilité des calculs.
\end{enumerate}
\end{procedure}

\section{Formules numériquement stables}

\begin{lecture}
Deux formules mathématiquement égales peuvent se comporter différemment sur un
ordinateur, dont les nombres ont une taille et une précision finies. Les formes
stables évitent de calculer des exponentielles immenses, puis de soustraire ou
de logarithmer des valeurs arrondies à $0$ ou $1$.
\end{lecture}

Évaluer $e^{-z}$ sans précaution peut déborder pour un score très négatif. Une
sigmoïde stable utilise deux branches :
\[
  \sigma(z)=
  \begin{cases}
    1/(1+e^{-z}),&z\geq0,\\
    e^z/(1+e^z),&z<0.
  \end{cases}
\]
Il est préférable de calculer directement la perte à partir du score,
\[
  \ell(y,z)=\log(1+e^z)-yz=\operatorname{softplus}(z)-yz,
\]
L'identité résulte directement du remplacement
$p=\sigma(z)$ dans la log-perte et constitue la forme usuelle de la perte
logistique \cite[sec.~7.1]{murphy2022}.
avec une primitive stable de type $\log(1+e^z)$, plutôt que de calculer
$\log(\sigma(z))$ après arrondi de la probabilité à $0$ ou $1$.

\begin{resultat}[Contrôle du gradient]
Avant tout entraînement, comparer le gradient analytique à une différence finie
centrée sur un cas test de faible dimension :
\[
  \frac{\partial J}{\partial\theta_j}\approx
  \frac{J(\thv+h e_j)-J(\thv-h e_j)}{2h}.
\]
Ce test détecte efficacement un signe inversé, un facteur $1/m$ oublié ou un
problème d'orientation des matrices. Il valide la dérivation, pas l'ensemble du
modèle.
\end{resultat}

\begin{figure}[htbp]
\centering
\begin{tikzpicture}
\begin{axis}[courseplot,axis lines=middle,xlabel={$\theta_j$},ylabel={$J$},
  xmin=-1.1,xmax=2.1,ymin=0,ymax=3.8,xtick={-.3,.2,.7},
  xticklabels={$\theta_j-h$,$\theta_j$,$\theta_j+h$},ytick=\empty,
  legend style={at={(.5,1.03)},anchor=south,legend columns=2,draw=none},samples=180]
  \addplot[accent,very thick,domain=-1.1:2.1]{(x-.7)^2+.45};
  \addlegendentry{$J(\theta_j)$}
  \addplot[warning,very thick,domain=-.3:.7]{-x+1.15};
  \addlegendentry{corde entre les deux évaluations}
  \addplot[only marks,mark=*,warning,mark size=2.8pt]
    coordinates {(-.3,1.45) (.7,.45)};
  \addplot[only marks,mark=*,ink,mark size=2.4pt] coordinates {(.2,.70)};
  \draw[dashed,gridgray] (axis cs:-.3,0)--(axis cs:-.3,1.45);
  \draw[dashed,gridgray] (axis cs:.7,0)--(axis cs:.7,.45);
  \node[anchor=west,align=left] at (axis cs:.82,1.15)
    {$h=0{,}5$\\pente $=(0{,}45-1{,}45)/(2h)=-1$};
\end{axis}
\end{tikzpicture}
\caption{Exemple numérique cohérent avec
$J(u)=(u-0{,}7)^2+0{,}45$, $\theta_j=0{,}2$ et $h=0{,}5$. Les deux points de la
corde appartiennent à la courbe. Sa pente vaut $-1$, exactement comme
$J'(0{,}2)=2(0{,}2-0{,}7)=-1$ dans cet exemple quadratique.}
\label{fig:finite-difference}
\end{figure}

\begin{procedure}{Évaluer les formules sans instabilité numérique}
\textbf{Entrée :} un score réel $z$, éventuellement de grande amplitude.
\textbf{Sorties :} une probabilité et une perte finies, calculées sans dépassement
de capacité.
\begin{enumerate}
  \item si $z\geq0$, utiliser $1/(1+e^{-z})$ ; sinon utiliser
        $e^z/(1+e^z)$, afin que l'exponentielle porte sur un nombre négatif ;
  \item calculer la perte depuis $z$ avec une primitive stable de softplus ;
  \item tester des scores extrêmes, par exemple $z=-1000,0,1000$, et vérifier
        l'absence de valeur infinie ou indéfinie ;
  \item sur un cas de faible dimension, comparer gradient analytique et différence
        finie avec une tolérance fixée avant l'entraînement complet.
\end{enumerate}
\end{procedure}

\chapter{Extension multi-classes un-contre-tous}
\label{ch:ovr}

\begin{lecture}
OvR transforme le problème à quatre maisons en quatre questions binaires :
« Gryffindor contre les autres », puis Hufflepuff, Ravenclaw et Slytherin. Pour
un nouvel élève, les quatre modèles produisent quatre scores ; la maison associée
au score maximal devient la prédiction.
\end{lecture}

Pour chaque classe $k\in\{1,\ldots,K\}$, on construit la cible binaire
\[
  y_{ik}=\ind[c_i=k].
\]
Cette réduction est la stratégie one-vs-rest : un classifieur binaire par classe,
la classe considérée étant positive et toutes les autres négatives
\cite[sec.~4.1]{rifkin2004}. On ajuste ensuite indépendamment un paramètre
$\thv_k$ pour chacun de ces problèmes. On obtient les scores
$z_k(\x)=\thv_k^\top\xt$ et les sorties binaires
$q_k(\x)=\sigma(z_k(\x))$. La règle de décision est
\[
  \widehat c(\x)=\Argmax_{k\in\{1,\ldots,K\}}q_k(\x)
  =\Argmax_{k\in\{1,\ldots,K\}}z_k(\x),
  \tag{3}\label{eq:ovr-decision}
\]
car $\sigma$ est strictement croissante et commune aux $K$ modèles.
La décision par maximum des réponses binaires est la règle OvR usuelle
\cite{rifkin2004}.

L'apprentissage produit ainsi $K$ vecteurs de paramètres
$\thv_1,\ldots,\thv_K$. Pour quatre maisons, une même matrice $\X$ est associée
successivement à quatre vecteurs cibles $\y_1,\ldots,\y_4$. À la prédiction, une
observation préparée est multipliée par chacun des quatre vecteurs de paramètres ;
les quatre scores obtenus sont comparés, et l'indice du maximum fournit la maison.

\begin{table}[htbp]
\centering
\begin{tabular}{lrrr}
\toprule
classe $k$ & score $z_k$ & sortie $q_k=\sigma(z_k)$ & rang \\
\midrule
Gryffindor & $-1{,}99$ & $0{,}12$ & $3$ \\
Hufflepuff & $-2{,}31$ & $0{,}09$ & $4$ \\
Ravenclaw  & $ 0{,}90$ & $0{,}71$ & $1$ \\
Slytherin  & $-1{,}21$ & $0{,}23$ & $2$ \\
\bottomrule
\end{tabular}
\caption{Décision OvR sur une observation. Ravenclaw maximise à la fois $z_k$ et
$q_k$, la même fonction strictement croissante $\sigma$ étant appliquée aux
quatre scores. La somme des sorties vaut ici $1{,}15$.}
\label{tab:ovr-example}
\end{table}

\begin{figure}[htbp]
\centering
\resizebox{\textwidth}{!}{%
\begin{tikzpicture}[x=.105cm,y=1cm]
  \node[anchor=east] at (0,3.5){G};
  \node[anchor=east] at (0,2.5){H};
  \node[anchor=east] at (0,1.5){R};
  \node[anchor=east] at (0,.5){S};
  \foreach \yy/\v in {3.5/12,2.5/9,1.5/71,.5/23}{
    \fill[accent!70] (0,\yy-.22) rectangle (\v,\yy+.22);
    \node[anchor=west] at (\v+1,\yy){\v\,\%};
  }
  \node[font=\bfseries] at (38,4.35){sorties OvR indépendantes : somme $115\,\%$};
  \begin{scope}[xshift=9.4cm]
    \node[anchor=east] at (0,3.5){G};
    \node[anchor=east] at (0,2.5){H};
    \node[anchor=east] at (0,1.5){R};
    \node[anchor=east] at (0,.5){S};
    \foreach \yy/\v in {3.5/5,2.5/3,1.5/82,.5/10}{
      \fill[warning!70] (0,\yy-.22) rectangle (\v,\yy+.22);
      \node[anchor=west] at (\v+1,\yy){\v\,\%};
    }
    \node[font=\bfseries] at (40,4.35){softmax des scores : somme $100\,\%$};
  \end{scope}
\end{tikzpicture}}
\caption{Comparaison sur les scores de la table~\ref{tab:ovr-example}. OvR
applique une sigmoïde séparée à chaque score ; softmax normalise conjointement
les exponentielles des scores. Les deux méthodes choisissent ici Ravenclaw, mais
leurs valeurs n'ont pas la même signification probabiliste.}
\label{fig:ovr-softmax}
\end{figure}

\begin{figure}[htbp]
  \centering
  \resizebox{\textwidth}{!}{%
  \begin{tikzpicture}[
    model/.style={draw=accent,fill=accentlight,rounded corners,minimum width=39mm,
      minimum height=8mm,align=center},
    score/.style={draw=ink!55,fill=white,rounded corners,minimum width=18mm,
      minimum height=8mm,align=center},arr/.style={-{Latex},draw=ink}]
    \node[model] (x) at (0,0) {même observation $\xt$};
    \node[model] (mG) at (5.2, 2.1) {Gryffindor contre le reste};
    \node[model] (mH) at (5.2,  .7) {Hufflepuff contre le reste};
    \node[model] (mR) at (5.2, -.7) {Ravenclaw contre le reste};
    \node[model] (mS) at (5.2,-2.1) {Slytherin contre le reste};
    \node[score] (sG) at (9.0, 2.1) {$z_G=-1{,}99$};
    \node[score] (sH) at (9.0,  .7) {$z_H=-2{,}31$};
    \node[score] (sR) at (9.0, -.7) {$z_R= 0{,}90$};
    \node[score] (sS) at (9.0,-2.1) {$z_S=-1{,}21$};
    \node[model,fill=warning!18,draw=warning] (max) at (12.6,0)
      {$\Argmax_k z_k=R$\\prédiction : Ravenclaw};
    \foreach \m/\s in {mG/sG,mH/sH,mR/sR,mS/sS}{
      \draw[arr] (x.east)--(\m.west);
      \draw[arr] (\m.east)--(\s.west);
      \draw[arr] (\s.east)--(max.west);
    }
  \end{tikzpicture}}
  \caption{Architecture OvR complète pour les quatre maisons. La même
  observation traverse les quatre modèles binaires, qui produisent quatre scores,
  ensuite comparés.}
\label{fig:ovr-architecture}
\end{figure}

\begin{procedure}{Ajuster et utiliser un modèle OvR}
\textbf{Entrées d'apprentissage :} $\X\in\R^{m\times(n+1)}$ et les classes
$c_i$. \textbf{État produit :} $K$ vecteurs $\thv_k$ associés à $K$ libellés.
\textbf{Sortie de prédiction :} un libellé de classe.
\begin{enumerate}
  \item pour la classe $k$, construire $y_{ik}=1$ si $c_i=k$, sinon $0$ ;
  \item ajuster $\thv_k$ avec $\X$ et ce vecteur cible, puis répéter pour les $K$ classes ;
  \item stocker les paramètres dans un ordre explicite, par exemple
        G, H, R, S, et conserver cet ordre avec le modèle ;
  \item pour une nouvelle ligne $\xt$, calculer
        $(z_G,z_H,z_R,z_S)=(\thv_G^\top\xt,\ldots,\thv_S^\top\xt)$ ;
  \item prendre la position de la plus grande valeur ; par exemple
        $(-1{,}99,-2{,}31,0{,}90,-1{,}21)$ renvoie la troisième position, donc Ravenclaw ;
  \item vérifier que quatre scores ont été produits et qu'aucun n'est non fini.
\end{enumerate}
\end{procedure}

\begin{vigilance}[Scores, probabilités binaires et probabilités multiclasses]
$q_k$ estime la probabilité de la cible binaire « classe $k$ contre le reste »
dans le modèle $k$. Les $K$ modèles étant ajustés séparément, chaque $q_k$ décrit
sa propre expérience binaire et la somme $\sum_k q_k$ prend une valeur libre. Ces
nombres servent au classement de l'équation~\eqref{eq:ovr-decision} ; leur usage
comme vecteur de probabilités multiclasses \rsym{calibration}{calibrées} demande une recalibration
préalable \cite{zadrozny2002}. Une régression logistique multinomiale à fonction
softmax fournit directement une \rsym{categorielle}{loi catégorielle} lorsque cet objectif
probabiliste est central \cite{hastie2009}.
\end{vigilance}

\begin{resultat}[Règle de décision OvR]
La décision OvR tient en deux opérations : calculer les quatre scores $z_k$, puis
retenir $\Argmax_k z_k$. La section suivante présente softmax comme un modèle
probabiliste multiclasses distinct, avec sa propre estimation.
\end{resultat}

La fonction \rsym{softmax}{softmax} mentionnée dans la comparaison est
\[
  \Pp(C=k\mid\x)
  =\frac{e^{z_k}}{\sum_{r=1}^K e^{z_r}}.
\]
Son dénominateur $\sum_r e^{z_r}$ couple les $K$ classes et impose une somme
égale à $1$. Dans OvR, le dénominateur $1+e^{-z_k}$ dépend du seul score $z_k$.
Cette différence algébrique fonde la différence d'interprétation ; sur
l'exactitude, les deux méthodes se départagent au cas par cas.

Les négatifs d'un problème OvR réunissent les trois autres maisons : ils sont
plus nombreux et plus hétérogènes que les positifs. Chaque sous-modèle affiche
donc une exactitude binaire élevée par construction. Le jugement porte sur la
décision multiclasses finale et sur les performances par maison.

\begin{figure}[htbp]
\centering
\begin{tikzpicture}[scale=1.15]
  \fill[accent!18] (-4,-3) rectangle (0,3);
  \fill[warning!18] (0,-3) rectangle (4,0);
  \fill[blue!12] (0,0) rectangle (4,3);
  \fill[orange!14] (-4,-3) rectangle (0,0);
  \draw[-{Latex},thick] (-4.2,0)--(4.4,0) node[right]{$x_1$};
  \draw[-{Latex},thick] (0,-3.2)--(0,3.4) node[above]{$x_2$};
  \draw[ink,very thick] (-4,0)--(4,0);
  \draw[ink,very thick] (0,-3)--(0,3);
  \node[align=center] at (-2,1.5){Gryffindor\\$z_G$ maximal};
  \node[align=center] at (1.6,1.5){Ravenclaw\\$z_R$ maximal};
  \node[align=center] at (-2,-1.5){Hufflepuff\\$z_H$ maximal};
  \node[align=center] at (2,-1.5){Slytherin\\$z_S$ maximal};
  \fill[ink] (3.0,.65) circle (2.2pt);
  \node[draw=ink!55,fill=white,rounded corners=1pt,
        inner sep=3pt,align=center] (obs) at (3.15,2.45)
        {observation $\x$\\à classer};
  \draw[-{Latex},ink!70] (obs.south)--(3.02,.75);
\end{tikzpicture}
\caption{Régions de décision OvR dans un espace à deux caractéristiques. Chaque
point rejoint la région dont le score domine. Les frontières entre deux régions
vérifient $z_k(\x)=z_r(\x)$ et restent linéaires tant que les scores le sont.}
\label{fig:ovr-regions}
\end{figure}

La règle du maximum s'applique dans deux situations remarquables. Quand les
quatre scores sont négatifs, elle retient le moins négatif. Quand plusieurs
scores sont fortement positifs, elle retient le plus grand. Dans les deux cas
elle produit une maison. Une politique de rejet forme une décision
supplémentaire, définie et validée à part.

Exemple de rejet : laisser la maison vide si le meilleur score descend sous $-2$,
ou si l'écart entre les deux meilleurs scores reste sous $0{,}1$. Ces seuils
viennent du protocole et introduisent leur propre arbitrage : davantage de rejets
réduit les mauvaises attributions et laisse davantage d'élèves sans réponse. Ils
se choisissent sur validation d'après le coût respectif d'un rejet et d'une
erreur, puis s'évaluent sur test.

\begin{resultat}[État conservé par le modèle OvR]
Avec $K$ classes, l'apprentissage conserve $K$ vecteurs de paramètres et l'ordre
qui associe chaque vecteur à son libellé. Une prédiction produit un vecteur de
$K$ scores dans ce même ordre, dont l'indice du maximum se convertit en libellé.
Cette correspondance appartient au modèle entraîné.
\end{resultat}

\chapter{Évaluation et expérience reproductible}
\label{ch:evaluation}

\section{Matrice de confusion et mesures}

\begin{lecture}
L'exactitude compte la proportion totale de décisions correctes. La matrice de
confusion conserve davantage d'information : ses lignes indiquent les maisons
réelles et ses colonnes les maisons prédites. Elle montre ainsi quelles classes
sont bien reconnues et lesquelles sont confondues.
\end{lecture}

Pour un test de taille $m_{\mathrm{te}}$, l'exactitude est
\[
  \operatorname{accuracy}=\frac1{m_{\mathrm{te}}}
  \sum_{i\in\mathcal I_{\mathrm{te}}}\ind[\widehat c_i=c_i].
\]
Cette quantité est l'estimateur empirique de la probabilité de classification
correcte sur un échantillon test indépendant \cite[chap.~7]{hastie2009}.
En français statistique, \emph{exactitude} évite la confusion avec la
\emph{précision} au sens anglophone de $\mathrm{precision}=TP/(TP+FP)$.

L'évaluation reçoit deux vecteurs de même longueur : les classes réelles
$(c_i)$ et les classes prédites $(\widehat c_i)$ du test. Chaque paire
$(c_i,\widehat c_i)$ incrémente une case d'une matrice $K\times K$. L'exactitude
est calculée depuis la diagonale ; les mesures par classe utilisent ensuite les
sommes des lignes et des colonnes. Aucun réentraînement n'intervient dans ces
comptages.

La matrice de confusion $M$, définie par
$M_{ab}=\#\{i:c_i=a,\widehat c_i=b\}$, révèle quelles maisons sont
confondues. Pour chaque classe, on complète par rappel, précision et score $F_1$.
La moyenne macro donne le même poids à chaque maison ; la moyenne pondérée tient
compte de leurs effectifs. Une valeur unique ne remplace pas ces diagnostics.

Le score $F_1$ d'une classe est la moyenne harmonique de sa précision $P$ et de
son rappel $R$ :
\[
  F_1=\frac{2PR}{P+R}.
\]
Il n'est élevé que si les deux quantités le sont. Avec $P=0{,}80$ et $R=0{,}50$,
on obtient $F_1\approx0{,}615$, et non leur moyenne arithmétique $0{,}65$.

Supposons ensuite deux classes : la première contient $900$ observations et a un
$F_1$ de $0{,}95$ ; la seconde en contient $100$ et a un $F_1$ de $0{,}40$. La
moyenne macro vaut $(0{,}95+0{,}40)/2=0{,}675$. La moyenne pondérée vaut
$(900\times0{,}95+100\times0{,}40)/1000=0{,}895$. Le second nombre paraît bien
meilleur parce que la classe majoritaire domine le calcul. Rapporter les deux
indique si l'on privilégie la performance par individu ou l'équilibre entre
classes.

\begin{table}[htbp]
\centering
\setlength{\tabcolsep}{9pt}
\begin{tabular}{lrrrr|r}
\toprule
réelle $\backslash$ prédite & G & H & R & S & total \\
\midrule
G & \cellcolor{accent!35}$42$ & $3$ & $2$ & $3$ & $50$ \\
H & $4$ & \cellcolor{accent!35}$35$ & $6$ & $5$ & $50$ \\
R & $1$ & $4$ & \cellcolor{accent!35}$43$ & $2$ & $50$ \\
S & $5$ & $2$ & $3$ & \cellcolor{accent!35}$40$ & $50$ \\
\midrule
total prédit & $52$ & $44$ & $54$ & $50$ & $200$ \\
\bottomrule
\end{tabular}
\caption{Matrice de confusion fictive. Les bonnes décisions occupent la
diagonale. L'exactitude vaut $(42+35+43+40)/200=0{,}80$, mais la ligne Hufflepuff
révèle un rappel inférieur aux autres classes.}
\label{tab:confusion-example}
\end{table}

Dans cet exemple, le rappel de Hufflepuff vaut $35/50=0{,}70$, tandis que sa
précision vaut $35/44\approx0{,}80$. La première quantité répond à « quelle part
des vrais Hufflepuff est retrouvée ? » ; la seconde à « parmi les élèves prédits
Hufflepuff, quelle part l'est réellement ? ». Ces dénominateurs différents
expliquent pourquoi les deux mesures ne sont pas interchangeables.

\begin{figure}[htbp]
\centering
\begin{tikzpicture}[x=.12cm,y=1cm]
  \draw[-{Latex},thick] (0,0)--(105,0) node[right]{\%};
  \foreach \x in {0,20,40,60,80,100}{
    \draw (\x,.08)--(\x,-.08) node[below]{\x};
  }
  \foreach \yy/\name/\rec/\prec in
    {4/Gryffindor/84/80.8,3/Hufflepuff/70/79.5,2/Ravenclaw/86/79.6,1/Slytherin/80/80}{
    \node[anchor=east] at (-2,\yy){\name};
    \fill[accent!75] (0,\yy+.05) rectangle (\rec,\yy+.35);
    \fill[warning!70] (0,\yy-.35) rectangle (\prec,\yy-.05);
  }
  \node[accent,anchor=west] at (5,5){rappel};
  \node[warning,anchor=west] at (25,5){précision};
\end{tikzpicture}
\caption{Rappel et précision par classe calculés à partir de la
table~\ref{tab:confusion-example}. Une exactitude globale de $80\,\%$ masque
notamment le rappel plus faible de Hufflepuff.}
\label{fig:metrics-comparison}
\end{figure}

\begin{procedure}{Construire les mesures de classification}
\textbf{Entrées :} $m_{\mathrm{te}}$ classes réelles et autant de classes
prédites, codées dans le même ordre de $K$ libellés. \textbf{Sorties :} une
matrice de confusion $M$, l'exactitude et les mesures par classe.
\begin{enumerate}
  \item initialiser $M\in\mathbb{N}^{K\times K}$ à zéro ; les lignes représentent le
        réel et les colonnes la prédiction ;
  \item pour chaque paire $(c_i,\widehat c_i)$, incrémenter $M_{c_i,\widehat c_i}$ ;
  \item vérifier que la somme de toutes les cases vaut $m_{\mathrm{te}}$ ;
  \item calculer l'exactitude comme somme diagonale divisée par $m_{\mathrm{te}}$ ;
  \item pour la classe $k$, calculer rappel = diagonale/somme de la ligne et
        précision = diagonale/somme de la colonne, puis $F_1$ ;
  \item signaler tout dénominateur nul et rapporter les valeurs par classe.
\end{enumerate}
\end{procedure}

\section{Protocole minimal d'évaluation}

\begin{lecture}
Un score isolé est impossible à interpréter sans connaître la manière dont il a
été obtenu. Le protocole enregistre les données utilisées, leur préparation, les
paramètres d'entraînement et les mesures finales afin que la même expérience
puisse être reproduite et contrôlée.
\end{lecture}

Une expérience exploitable doit conserver :
\begin{enumerate}
  \item les variables retenues et la politique de valeurs manquantes ;
  \item la partition, sa stratification et sa graine ;
  \item les statistiques de préparation apprises sur le train ;
  \item $\alpha$, $\eps$, le nombre maximal d'itérations et, le cas échéant,
        $\lambda$ ;
  \item pour chaque classe, le coût initial, le coût final, le nombre
        d'itérations et la norme finale du gradient ;
  \item sur le test : matrice de confusion, exactitude et mesures par classe.
\end{enumerate}

Deux contrôles simples donnent du sens au résultat. Le modèle doit dépasser une
référence majoritaire qui prédit toujours la maison la plus fréquente. Puis une
permutation aléatoire des étiquettes doit ramener la performance vers le hasard ;
une performance anormalement élevée après permutation signale souvent une fuite
de données ou une erreur de protocole.

Exemple : si Gryffindor représente $35\,\%$ du test, prédire Gryffindor pour tout
le monde donne déjà $35\,\%$ d'exactitude. Un modèle obtenant $38\,\%$ ne doit
donc pas être comparé seulement au hasard uniforme de $25\,\%$, mais aussi à
cette référence majoritaire. Pour le contrôle par permutation, on mélange les
maisons entre les lignes avant l'entraînement tout en conservant les notes. Le
lien entre caractéristiques et réponse est alors détruit. Sur quatre classes
équilibrées, une exactitude proche de $25\,\%$ est attendue ; obtenir encore
$80\,\%$ indiquerait presque certainement que la réponse ou une information qui
la révèle s'est introduite dans les variables ou dans la préparation.

\begin{procedure}{Exécuter une évaluation contrôlée}
\textbf{Entrées :} modèle figé, statistiques de préparation, indices de test et
classes réelles. \textbf{Sortie :} un compte rendu reproductible, sans
modification du modèle.
\begin{enumerate}
  \item sélectionner les lignes désignées par $\mathcal I_{\mathrm{te}}$ et leur
        appliquer les statistiques d'apprentissage déjà conservées ;
  \item produire exactement une prédiction par ligne, sans mise à jour des paramètres ;
  \item vérifier l'égalité des longueurs entre prédictions et classes réelles ;
  \item construire matrice de confusion, exactitude et mesures par classe ;
  \item comparer l'exactitude à celle de la classe majoritaire et exécuter
        séparément le contrôle par permutation sur une expérience d'apprentissage ;
  \item enregistrer résultats, paramètres, hyperparamètres, graine et indices.
\end{enumerate}
\end{procedure}

\section{Validation croisée : plis et répétitions}

\begin{lecture}
Une seule partition peut avantager ou désavantager un modèle par hasard. La
validation croisée fait tourner le rôle de validation entre plusieurs blocs :
chaque observation est évaluée une fois tout en restant hors de l'apprentissage
correspondant. Répéter l'ensemble du découpage mesure la sensibilité à cette
répartition aléatoire.
\end{lecture}

Un \emph{pli} désigne un bloc d'observations. Dans une validation croisée
à cinq plis, on découpe les données disponibles pour la sélection du modèle en
cinq blocs disjoints $B_1,\ldots,B_5$. On entraîne cinq fois le modèle : à chaque
fois, quatre blocs servent à l'ajustement et le bloc restant sert à la validation.

\begin{figure}[htbp]
\centering
\begingroup
\setlength{\tabcolsep}{4pt}
\renewcommand{\arraystretch}{1.35}
\begin{tabular}{r*{5}{>{\centering\arraybackslash}p{13mm}}}
& $B_1$ & $B_2$ & $B_3$ & $B_4$ & $B_5$ \\
validation sur $B_1$ & \cellcolor{warning!70}V & \cellcolor{accent!55}A & \cellcolor{accent!55}A & \cellcolor{accent!55}A & \cellcolor{accent!55}A \\
validation sur $B_2$ & \cellcolor{accent!55}A & \cellcolor{warning!70}V & \cellcolor{accent!55}A & \cellcolor{accent!55}A & \cellcolor{accent!55}A \\
validation sur $B_3$ & \cellcolor{accent!55}A & \cellcolor{accent!55}A & \cellcolor{warning!70}V & \cellcolor{accent!55}A & \cellcolor{accent!55}A \\
validation sur $B_4$ & \cellcolor{accent!55}A & \cellcolor{accent!55}A & \cellcolor{accent!55}A & \cellcolor{warning!70}V & \cellcolor{accent!55}A \\
validation sur $B_5$ & \cellcolor{accent!55}A & \cellcolor{accent!55}A & \cellcolor{accent!55}A & \cellcolor{accent!55}A & \cellcolor{warning!70}V \\
\end{tabular}

\vspace{.6em}
\textcolor{accent}{A : apprentissage}\qquad
\textcolor{warning}{V : validation}
\endgroup
\caption{Validation croisée à cinq plis. Chaque observation sert quatre fois à
l'apprentissage et exactement une fois à la validation. Le test final reste à
l'écart de l'ensemble de ce mécanisme.}
\label{fig:five-fold}
\end{figure}

Supposons que les exactitudes de validation des cinq tours soient
$0{,}79$, $0{,}82$, $0{,}77$, $0{,}81$ et $0{,}80$. Leur moyenne vaut
$0{,}798$. Leur dispersion montre aussi que le résultat dépend quelque peu des
observations placées dans le bloc de validation ; annoncer seulement « environ
$0{,}80$ » masque cette information.

Une \emph{répétition} consiste à refaire tout le découpage en cinq blocs avec une
autre permutation aléatoire, puis à recommencer les cinq entraînements. Les
observations ne changent pas ; seule leur affectation aux blocs change. Deux
répétitions produisent ici dix mesures :

\begin{table}[htbp]
\centering
\begin{tabular}{lccccc|c}
\toprule
& pli 1 & pli 2 & pli 3 & pli 4 & pli 5 & moyenne \\
\midrule
répétition 1 & $0{,}79$ & $0{,}82$ & $0{,}77$ & $0{,}81$ & $0{,}80$ & $0{,}798$ \\
répétition 2 & $0{,}81$ & $0{,}78$ & $0{,}80$ & $0{,}76$ & $0{,}83$ & $0{,}796$ \\
\midrule
\multicolumn{6}{r}{moyenne des dix validations} & $0{,}797$ \\
\bottomrule
\end{tabular}
\caption{Exemple fictif de validation croisée répétée. Les dix nombres ne sont
pas dix évaluations finales indépendantes : les ensembles d'apprentissage se
recouvrent fortement. Ils décrivent néanmoins la sensibilité au découpage.}
\label{tab:repeated-cv}
\end{table}

\begin{procedure}{Comparer des configurations par validation croisée}
\textbf{Entrées :} les données d'apprentissage, une liste de configurations et
un nombre de plis $q$. Un pli est une liste d'indices utilisée une fois pour la
validation. \textbf{Sortie :} une configuration choisie sans consulter le test.
\begin{enumerate}
  \item répartir les indices en $q$ blocs disjoints et enregistrer ces blocs ;
  \item pour une configuration donnée et un bloc $B_r$, ajuster sur tous les
        indices hors de $B_r$, puis mesurer sur $B_r$ ;
  \item répéter pour $r=1,\ldots,q$ : chaque ligne sert ainsi une fois à la validation ;
  \item calculer la moyenne et la dispersion des $q$ mesures ;
  \item recommencer avec exactement les mêmes blocs pour chaque configuration ;
  \item choisir selon une règle fixée, puis seulement exécuter le protocole final et le test.
\end{enumerate}
\end{procedure}

Pour comparer deux modèles A et B, on utilise exactement les mêmes blocs pour les
deux. Si, sur un pli donné, A obtient $0{,}81$ et B $0{,}79$, la différence est
$+0{,}02$. Répéter ce calcul pli par pli sépare mieux l'effet du modèle de la
difficulté particulière d'un bloc. Employer des découpages différents ajouterait
inutilement la variation des données à la comparaison.

\begin{vigilance}[Incertitude de l'évaluation]
Une exactitude test est elle-même une estimation soumise à la variabilité
d'échantillonnage. Sur $200$ observations, une différence de deux bonnes
prédictions ne représente qu'un point de pourcentage et ne suffit généralement
pas à conclure qu'un modèle est réellement supérieur. La comparaison de modèles
doit conserver les mêmes partitions. Les plis et répétitions de la section
précédente mesurent la sensibilité pendant la sélection ; ils ne transforment pas
le test final en certitude.
\end{vigilance}

Par exemple, A classe correctement $162$ élèves sur $200$ et B en classe $160$.
L'écart d'exactitude vaut $0{,}01$, mais il faut regarder les décisions appariées.
Si A corrige deux erreurs de B sans introduire aucune nouvelle erreur, l'écart est
cohérent sur ces deux cas. Si A et B se trompent chacun sur vingt élèves presque
tous différents, le même total cache des comportements très différents. La
matrice de confusion et la liste des désaccords sont alors plus informatives que
le seul écart d'un point.

\begin{resultat}[Résultats produits par l'évaluation]
L'évaluation finale produit le vecteur des classes prédites, la matrice de
confusion, l'exactitude globale et les mesures par classe. Les historiques de
coût et de norme du gradient décrivent la convergence de l'entraînement ; ils ne
remplacent pas les mesures calculées sur le test.
\end{resultat}

\chapter{Entraînement et prédiction}
\label{ch:implementation}

\begin{lecture}
Deux programmes suffisent à produire le classifieur. Le premier lit les élèves
étiquetés et écrit les paramètres ; le second lit ces paramètres et attribue une
maison à chaque nouvel élève. Ce chapitre donne les colonnes qui entrent dans le
calcul, les dimensions de chaque tableau, les valeurs des hyperparamètres et le
contenu des fichiers échangés. Les nombres cités proviennent de
\texttt{dataset\_train.csv}.
\end{lecture}

\section{Colonnes du fichier}
\label{sec:features}

Le fichier d'apprentissage contient $1\,600$ lignes et dix-neuf colonnes.

\begin{table}[htbp]
\centering
\begin{tabularx}{\textwidth}{>{\ttfamily\footnotesize}l X}
\toprule
\textrm{colonne} & \textrm{contenu et emploi} \\
\midrule
Index & entier de $0$ à $1\,599$. Recopié dans \texttt{houses.csv} pour apparier chaque prédiction à son élève \\
Hogwarts House & la réponse $c_i$. Quatre valeurs, d'effectifs $529$, $443$, $327$ et $301$. Source des cibles binaires $y_{ik}$ \\
First Name, Last Name & identité de l'élève. Écartées \\
Birthday & date de naissance au format \texttt{AAAA-MM-JJ}, de $1996$ à $2001$. Écartée d'après la section~\ref{sec:hand-birthday} \\
Best Hand & \texttt{Left} pour $790$ élèves, \texttt{Right} pour $810$. Écartée d'après la section~\ref{sec:hand-birthday} \\
\textrm{treize colonnes de notes} & nombres réels. Variables candidates \\
\bottomrule
\end{tabularx}
\caption{Composition du fichier d'apprentissage.}
\label{tab:columns}
\end{table}

\section{Best Hand et Birthday}
\label{sec:hand-birthday}

Ces deux colonnes sont remplies sur les $1\,600$ lignes et demandent un codage
numérique avant d'entrer dans le modèle. Trois mesures décident de leur sort.

La proportion de gauchers varie de $48{,}1\,\%$ chez Ravenclaw à $51{,}4\,\%$
chez Gryffindor, pour $49{,}4\,\%$ sur l'ensemble. Le tableau de contingence
$4\times2$ donne un \rsym{chi2}{$\chi^2$} de $0{,}94$ à trois degrés de liberté, quand le
seuil à $5\,\%$ vaut $7{,}81$.

L'écart entre la note moyenne des gauchers et celle des droitiers reste sous
$0{,}10$ écart-type pour les treize cours, avec un maximum de $0{,}095$ pour
Defense Against the Dark Arts. L'\rsym{erreurtype}{erreur type} d'un tel écart vaut
$\sqrt{1/790+1/810}=0{,}050$ écart-type : les treize valeurs tiennent à moins de
deux erreurs types de zéro. Les maisons, elles, se séparent de $1{,}86$ à
$2{,}66$ écarts-types sur ces mêmes cours.

Le mois de naissance donne un $\chi^2$ de $33{,}4$ à $33$ degrés de liberté.
L'année donne $28{,}1$ à $15$ degrés de liberté, au-dessus du seuil à $5\,\%$ qui
vaut $25{,}0$ : une association faible, portée en partie par les $40$ élèves nés
en $1996$ contre environ $320$ pour chacune des années suivantes.

Le test décisif porte sur l'exactitude. Best Hand codée $\pm1$, puis l'année de
naissance standardisée, ajoutées chacune comme onzième variable, laissent
l'exactitude hors échantillon à $0{,}9832$ sur cinq graines de découpage,
identique aux quatre décimales. Les coefficients de Best Hand s'échelonnent de
$-0{,}20$ à $+0{,}19$, contre $2{,}39$ pour Astronomy chez Hufflepuff. Les deux
colonnes restent hors du modèle.

\section{Redondance}

\begin{lecture}
Deux colonnes qui portent la même information gênent le calcul. Le modèle peut
augmenter le coefficient de l'une et diminuer celui de l'autre sans rien changer
à ses prédictions : ses coefficients dérivent alors sans jamais se fixer. Une des
deux colonnes suffit.
\end{lecture}

Le jeu de données vérifie
\[
  \text{Astronomy}=-100\times\text{Defense Against the Dark Arts}
\]
sur les $1\,537$ lignes où les deux valeurs sont présentes, avec un écart maximal
de $4{,}5\cdot10^{-13}$ et une \rsym{correlation}{corrélation} de $-1$ à six décimales. Le
chapitre~\ref{ch:likelihood} établit la conséquence : deux colonnes
\rsym{colinearite}{proportionnelles} rendent \rsym{xtwx}{$\X^\top W\X$} singulière et laissent la fonction de coût
plate le long d'une direction. Les coefficients dérivent alors librement sur
cette direction pendant que les scores restent identiques. Astronomy est
conservée, Defense Against the Dark Arts sert à l'imputation décrite plus bas.

\section{Pouvoir discriminant}

\begin{lecture}
Une note aide à reconnaître une maison lorsque les quatre maisons y obtiennent
des résultats différents. Si les quatre ont la même moyenne, la note ne distingue
personne. Le nombre $\delta_j$ ci-dessous mesure cet écart entre maisons, dans
une unité qui permet de comparer des cours notés sur des échelles très
différentes.
\end{lecture}

Pour la variable $j$, l'étendue des moyennes par maison rapportée à l'écart-type
global vaut
\[
  \delta_j=\frac{\max_k\bar r_{jk}-\min_k\bar r_{jk}}{s_j}.
\]
Un $\delta_j$ de $2$ signifie que les maisons extrêmes ont des moyennes séparées
par deux écarts-types : leurs distributions se chevauchent peu et la variable
sépare. Un $\delta_j$ de $0{,}07$ place les quatre moyennes à l'intérieur d'un
quatorzième d'écart-type : les quatre histogrammes se superposent.

\begin{table}[htbp]
\centering
\setlength{\tabcolsep}{7pt}
\begin{tabular}{lrrl}
\toprule
variable & $\delta_j$ & manquantes & décision \\
\midrule
\cellcolor{warning!14}Arithmancy                & \cellcolor{warning!14}$0{,}068$ & \cellcolor{warning!14}$34$ & \cellcolor{warning!14}écartée : homogène \\
\cellcolor{warning!14}Care of Magical Creatures & \cellcolor{warning!14}$0{,}147$ & \cellcolor{warning!14}$40$ & \cellcolor{warning!14}écartée : homogène \\
\midrule
Ancient Runes                & $1{,}860$ & $35$ & retenue \\
Herbology                    & $1{,}879$ & $33$ & retenue \\
\cellcolor{warning!14}Defense Against the Dark Arts & \cellcolor{warning!14}$1{,}909$ & \cellcolor{warning!14}$31$ & \cellcolor{warning!14}écartée : redondante \\
Astronomy                    & $1{,}911$ & $32$ & retenue \\
Muggle Studies               & $2{,}038$ & $35$ & retenue \\
Potions                      & $2{,}076$ & $30$ & retenue \\
History of Magic             & $2{,}221$ & $43$ & retenue \\
Transfiguration              & $2{,}284$ & $34$ & retenue \\
Divination                   & $2{,}368$ & $39$ & retenue \\
Charms                       & $2{,}466$ & $0$  & retenue \\
Flying                       & $2{,}657$ & $0$  & retenue \\
\bottomrule
\end{tabular}
\caption{Les treize notes classées par pouvoir discriminant croissant. Arithmancy
et Care of Magical Creatures sont en dessous des autres d'un facteur douze ; la
coupure tombe dans un intervalle vide.}
\label{tab:feature-selection}
\end{table}

\begin{resultat}[Variables retenues]
Astronomy, Herbology, Divination, Muggle Studies, Ancient Runes,
History of Magic, Transfiguration, Potions, Charms, Flying.
Donc $n=10$, $\X\in\R^{m\times11}$ et chaque $\thv_k\in\R^{11}$.
\end{resultat}

\section{Valeurs manquantes}

Onze des treize colonnes de notes comportent des cases vides, entre $30$ et $43$
chacune, réparties sur $349$ lignes distinctes.

Les $1\,600$ lignes possèdent au moins une des deux valeurs Astronomy ou Defense
Against the Dark Arts. La relation de proportionnalité donne donc les $32$
valeurs manquantes d'Astronomy à partir de la colonne écartée, au facteur $-100$
près. Ces $32$ cases reçoivent leur valeur exacte, et Astronomy devient complète.

Restent alors $239$ lignes portant au moins une case vide parmi les dix variables
retenues, soit $14{,}9\,\%$ de l'échantillon : les supprimer ramènerait
l'apprentissage à $1\,361$ élèves. L'imputation médiane du
chapitre~\ref{ch:data} les conserve toutes, et porte sur les sept colonnes encore
incomplètes. Charms, Flying et Astronomy sont pleines.

\section{Matrice de design}

Le programme applique dans l'ordre : reconstruction d'Astronomy, imputation
médiane, standardisation, ajout de la colonne de biais.

\begin{procedure}{Construire $\X$}
\textbf{Entrée :} un fichier de données. \textbf{Sorties :} $\X$, le vecteur des
identifiants \texttt{Index}, et le vecteur des classes lorsque la colonne est
remplie.
\begin{enumerate}
  \item lire les lignes dans l'ordre du fichier ; cet ordre est celui de
        \texttt{houses.csv} ;
  \item extraire les dix colonnes retenues, dans l'ordre fixé par le fichier de
        poids ;
  \item convertir les cases en nombres réels ; une case vide devient une valeur
        manquante ;
  \item remplir les valeurs manquantes d'Astronomy par $-100\,r_{\mathrm{DADA}}$,
        puis les valeurs manquantes restantes par la \rsym{mediane}{médiane} d'apprentissage de
        leur colonne ;
  \item appliquer $x_{ij}=(r_{ij}-\mu_j)/s_j$ avec les $\mu_j$ et $s_j$
        d'apprentissage ;
  \item préfixer chaque ligne d'une coordonnée égale à $1$, ce qui donne
        $\X\in\R^{m\times11}$.
\end{enumerate}
\end{procedure}

Sur \texttt{dataset\_train.csv}, $m=1\,600$ et $\X$ mesure $1\,600\times11$. Sur
\texttt{dataset\_test.csv}, $m=400$ et $\X$ mesure $400\times11$. Le découpage
stratifié $80\,\%/20\,\%$ du chapitre~\ref{ch:evaluation} partage les $1\,600$
lignes d'apprentissage en $1\,278$ et $322$.

\section{Le programme d'entraînement}

\begin{lecture}
L'entraînement répète quatre calculs jusqu'à ce que les coefficients cessent de
bouger : un score pour chaque élève, la probabilité qui lui correspond, l'écart
entre cette probabilité et la bonne réponse, puis une correction des coefficients
proportionnelle à cet écart. La même boucle tourne quatre fois, une par maison.
\end{lecture}

L'entraînement traite les quatre maisons l'une après l'autre. Pour la maison $k$,
il forme la cible $\y_k\in\{0,1\}^{1600}$ par $y_{ik}=\ind[c_i=k]$, part de
$\thv_k=\mathbf 0$, puis répète quatre opérations :
\[
  \mathbf z=\X\thv_k\;(m),\qquad
  \mathbf p=\sigma(\mathbf z)\;(m),\qquad
  \mathbf g=\tfrac1m\X^\top(\mathbf p-\y_k)\;(11),\qquad
  \thv_k\leftarrow\thv_k-\alpha\mathbf g.
\]
Le coût $J$ de l'équation~\eqref{eq:loss} sert au critère d'arrêt, le gradient
$\mathbf g$ est l'équation~\eqref{eq:gradient}. Le départ à $\thv_k=\mathbf 0$
donne un résultat reproductible d'une exécution à l'autre, la convexité établie
au chapitre~\ref{ch:likelihood} garantissant l'unicité de la limite à direction
plate près.

\begin{procedure}{logreg\_train}
\textbf{Entrée :} \texttt{dataset\_train.csv}. \textbf{Sortie :} le fichier de
poids.
\begin{enumerate}
  \item calculer médianes, $\mu_j$ et $s_j$ sur les lignes d'apprentissage, puis
        construire $\X$ ;
  \item pour $k=1,\ldots,4$, former $\y_k$ et itérer les quatre opérations
        ci-dessus jusqu'à $|J^{(t+1)}-J^{(t)}|/\max(1,|J^{(t)}|)<\eps$ ou
        épuisement du plafond d'itérations ;
  \item relever pour chaque maison le coût initial, le coût final, le nombre
        d'itérations et $\|\mathbf g\|_2$ ;
  \item écrire le fichier de poids avec les statistiques de préparation.
\end{enumerate}
\end{procedure}

\begin{table}[htbp]
\centering
\begin{tabular}{rrrrr}
\toprule
$\alpha$ & $\eps$ & itérations & $\|\nabla J\|_2$ final & exactitude hors échantillon \\
\midrule
$0{,}5$ & $10^{-6}$ & $1\,817$ & $1{,}4\cdot10^{-3}$ & $0{,}981$ \\
$1$     & $10^{-6}$ & $1\,461$ & $1{,}0\cdot10^{-3}$ & $0{,}981$ \\
$2$     & $10^{-6}$ & $1\,150$ & $7{,}1\cdot10^{-4}$ & $0{,}981$ \\
$1$     & $10^{-7}$ & $5\,448$ & $3{,}2\cdot10^{-4}$ & $0{,}981$ \\
\bottomrule
\end{tabular}
\caption{Itérations consommées par la maison la plus lente, sur les $1\,278$
lignes d'apprentissage du découpage stratifié. Diviser $\eps$ par dix multiplie
les itérations par près de quatre et laisse les $322$ décisions du test interne
inchangées : la convergence numérique se poursuit après que le classement s'est
stabilisé.}
\label{tab:training-defaults}
\end{table}

\begin{resultat}[Configuration de départ]
$\alpha=1$, $\eps=10^{-6}$, plafond de $6\,000$ itérations, $\thv_k=\mathbf 0$.
Cette configuration atteint $0{,}981$ sur un test interne de $322$ élèves, et
$0{,}983$ en moyenne sur cinq graines de découpage.
\end{resultat}

\section{Le fichier de poids}
\label{sec:weights}

La prédiction a besoin de trois choses : l'ordre des dix variables, les
statistiques qui les transforment, et les quatre vecteurs de paramètres associés
à leurs libellés. Les statistiques figurent dans le fichier parce qu'elles sont
estimées sur l'apprentissage ; les recalculer sur \texttt{dataset\_test.csv}
placerait les scores dans une autre échelle et ferait chuter l'exactitude.

\begin{fichier}[weights.csv]
feature,median,mu,sigma\\
Astronomy,258.1,38.6,520.8\\
Herbology,3.5,1.2,5.2\\
{[}...{]}\\
Flying,-2.5,22.0,97.6\\
class,bias,Astronomy,Herbology,...,Flying\\
Gryffindor,-3.54,0.50,-1.81,...,1.61\\
Hufflepuff,-2.36,2.48,1.63,...,-0.48\\
Ravenclaw,-2.56,-1.37,0.74,...,-0.17\\
Slytherin,-3.77,-2.02,-0.70,...,-0.59
\end{fichier}

Ces valeurs sont celles produites par le programme de la
section~\ref{sec:pseudo-train} sur les $1\,600$ lignes du fichier
d'apprentissage. L'écart-type d'Astronomy vaut $521$ et celui d'Herbology
$5{,}2$ : un facteur cent entre deux colonnes que le même $\alpha$ doit faire
progresser, ce que la standardisation ramène à $1$.

L'ordre des colonnes et l'ordre des classes appartiennent au modèle. Un vecteur
$\thv_k$ s'applique aux variables dans l'ordre où il a été ajusté, et l'indice
renvoyé par $\Argmax$ se lit dans la liste des libellés enregistrée à côté de
lui. Les deux ordres figurent explicitement dans le fichier.

\section{Primitives à écrire}

Le sujet écarte les fonctions toutes faites de statistique descriptive. Cinq
primitives suffisent, toutes en une passe ou deux sur un tableau.

\begin{pseudo}[primitives]
fonction moyenne(v) :\\
\hspace*{2em}s $\leftarrow$ 0\\
\hspace*{2em}pour x dans v : s $\leftarrow$ s + x\\
\hspace*{2em}retourner s / longueur(v)\\
\\
fonction ecart\_type(v) :\\
\hspace*{2em}m $\leftarrow$ moyenne(v) ; s $\leftarrow$ 0\\
\hspace*{2em}pour x dans v : s $\leftarrow$ s + (x - m)$^2$\\
\hspace*{2em}retourner racine(s / longueur(v))\\
\\
fonction mediane(v) :\\
\hspace*{2em}t $\leftarrow$ trier(v) ; n $\leftarrow$ longueur(t)\\
\hspace*{2em}si n impair : retourner t[n // 2]\\
\hspace*{2em}retourner (t[n // 2 - 1] + t[n // 2]) / 2\\
\\
fonction sigmoide(z) :\\
\hspace*{2em}si z $\geq$ 0 : retourner 1 / (1 + exp(-z))\\
\hspace*{2em}retourner exp(z) / (1 + exp(z))\\
\\
fonction softplus(z) :\\
\hspace*{2em}retourner max(z, 0) + log(1 + exp(-abs(z)))
\end{pseudo}

Les deux dernières viennent du chapitre~\ref{ch:optim} : leur exponentielle porte
toujours sur un nombre négatif ou nul, ce qui les garde finies pour
$z=\pm1000$.

\section{Pseudo-code de logreg\_train}
\label{sec:pseudo-train}

Le tableau \texttt{X} a $m$ lignes et $11$ colonnes, \texttt{theta[k]} $11$
composantes, \texttt{y} $m$ composantes. Les indices de colonne vont de $0$ à
$10$, la colonne $0$ portant le biais.

\begin{pseudo}[logreg\_train dataset\_train.csv]
VARIABLES $\leftarrow$ [Astronomy, Herbology, Divination, Muggle Studies,\\
\hspace*{5em}Ancient Runes, History of Magic, Transfiguration,\\
\hspace*{5em}Potions, Charms, Flying]\\
MAISONS $\leftarrow$ [Gryffindor, Hufflepuff, Ravenclaw, Slytherin]\\
ALPHA $\leftarrow$ 1.0 ; EPS $\leftarrow$ 1e-6 ; MAX\_ITER $\leftarrow$ 6000\\
\\
lignes $\leftarrow$ lire\_csv(argv[1])\\
\\
pour i dans lignes :\\
\hspace*{2em}si lignes[i][Astronomy] est vide et lignes[i][DADA] est rempli :\\
\hspace*{4em}lignes[i][Astronomy] $\leftarrow$ -100 * lignes[i][DADA]\\
\\
pour f dans VARIABLES :\\
\hspace*{2em}presentes $\leftarrow$ [lignes[i][f] pour i tel que lignes[i][f] rempli]\\
\hspace*{2em}med[f] $\leftarrow$ mediane(presentes)\\
\hspace*{2em}pour i dans lignes :\\
\hspace*{4em}si lignes[i][f] est vide : lignes[i][f] $\leftarrow$ med[f]\\
\hspace*{2em}mu[f] $\leftarrow$ moyenne(colonne f)\\
\hspace*{2em}sd[f] $\leftarrow$ ecart\_type(colonne f)\\
\hspace*{2em}si sd[f] = 0 : erreur("colonne constante", f)\\
\\
X $\leftarrow$ tableau[m][11]\\
pour i de 0 a m-1 :\\
\hspace*{2em}X[i][0] $\leftarrow$ 1.0\\
\hspace*{2em}pour j de 0 a 9 :\\
\hspace*{4em}f $\leftarrow$ VARIABLES[j]\\
\hspace*{4em}X[i][j+1] $\leftarrow$ (lignes[i][f] - mu[f]) / sd[f]\\
\\
pour k de 0 a 3 :\\
\hspace*{2em}y $\leftarrow$ [1 si lignes[i][House] = MAISONS[k] sinon 0]\\
\hspace*{2em}th $\leftarrow$ [0.0] * 11\\
\hspace*{2em}J\_prec $\leftarrow$ indefini\\
\hspace*{2em}pour t de 1 a MAX\_ITER :\\
\hspace*{4em}J $\leftarrow$ 0 ; g $\leftarrow$ [0.0] * 11\\
\hspace*{4em}pour i de 0 a m-1 :\\
\hspace*{6em}z $\leftarrow$ 0\\
\hspace*{6em}pour j de 0 a 10 : z $\leftarrow$ z + th[j] * X[i][j]\\
\hspace*{6em}J $\leftarrow$ J + softplus(z) - y[i] * z\\
\hspace*{6em}e $\leftarrow$ sigmoide(z) - y[i]\\
\hspace*{6em}pour j de 0 a 10 : g[j] $\leftarrow$ g[j] + e * X[i][j]\\
\hspace*{4em}J $\leftarrow$ J / m\\
\hspace*{4em}pour j de 0 a 10 : g[j] $\leftarrow$ g[j] / m\\
\hspace*{4em}si J\_prec defini et abs(J\_prec - J) / max(1, abs(J)) < EPS :\\
\hspace*{6em}sortir de la boucle\\
\hspace*{4em}J\_prec $\leftarrow$ J\\
\hspace*{4em}pour j de 0 a 10 : th[j] $\leftarrow$ th[j] - ALPHA * g[j]\\
\hspace*{2em}theta[k] $\leftarrow$ th\\
\\
ecrire\_poids("weights.csv", VARIABLES, med, mu, sd, MAISONS, theta)
\end{pseudo}

La boucle sur $j$ dans le calcul de \texttt{z} et de \texttt{g} correspond au
produit $\X\thv_k$ et au produit $\X^\top(\mathbf p-\y_k)$. Une bibliothèque de
calcul matriciel remplace les deux boucles internes par
\texttt{z = X @ th} et \texttt{g = X.T @ (p - y) / m}, avec le même résultat.

L'ordre des opérations compte sur deux points. La division de \texttt{g} par $m$
arrive avant la mise à jour, et \texttt{th[j]} est modifié après le calcul
complet de \texttt{g} : les onze composantes descendent depuis le même vecteur.

Exécuté sur les $1\,600$ lignes, ce programme consomme $436$ itérations pour
Hufflepuff, $460$ pour Ravenclaw, $836$ pour Slytherin et $1\,069$ pour
Gryffindor, sous le plafond de $6\,000$. Les coûts finaux vont de $0{,}044$ à
$0{,}069$, et le modèle reclasse correctement $0{,}982$ de son propre ensemble
d'apprentissage.

\section{Le programme de prédiction}

La prédiction lit les paramètres, construit $\X$ avec les statistiques du fichier
de poids, calcule quatre scores par ligne et retient le plus grand, selon
l'équation~\eqref{eq:ovr-decision}. La sigmoïde est croissante et commune aux
quatre modèles : comparer les $z_k$ donne le même classement que comparer les
$q_k$.

La remarque s'étend à \rsym{softmax}{softmax}, dont le dénominateur est commun aux
quatre classes : la section~\ref{sec:lois} en donne l'argument. Les trois
quantités $z_k$, $\sigma(z_k)$ et $\mathrm{softmax}(z)_k$ désignent le même
argmax, pour tout jeu de scores. Sur les $1\,600$ lignes du fichier
d'apprentissage, les trois règles produisent la même maison, avec zéro ex æquo.

Deux mesures situent la difficulté du problème. L'écart entre le meilleur score
et le second vaut $2{,}05$ au minimum et $8{,}69$ en médiane. Le maximum des
quatre scores est positif sur chaque ligne, et une seule classe y dépasse $2$.
Les quatre maisons occupent donc des régions franchement séparées, ce qui éclaire
le résultat suivant : une régression softmax multinomiale, ajustée conjointement
sur les mêmes découpages, produit les $322$ mêmes prédictions que OvR à chaque
graine, pour une exactitude moyenne identique de $0{,}983$. Sur ce jeu de
données, le choix entre les deux modèles porte sur l'interprétation des sorties,
traitée au chapitre~\ref{ch:ovr}.

\begin{procedure}{logreg\_predict}
\textbf{Entrées :} \texttt{dataset\_test.csv} et le fichier de poids.
\textbf{Sortie :} \texttt{houses.csv}.
\begin{enumerate}
  \item lire l'ordre des variables, les statistiques de préparation, l'ordre des
        classes et les quatre $\thv_k$ ;
  \item construire $\X\in\R^{400\times11}$ avec ces statistiques ;
  \item calculer $\mathbf Z=\X\Theta^\top\in\R^{400\times4}$, où $\Theta$ empile
        les quatre vecteurs de paramètres ;
  \item pour chaque ligne, prendre l'indice du maximum et le convertir en
        libellé ;
  \item écrire l'en-tête puis $400$ lignes, dans l'ordre du fichier d'entrée.
\end{enumerate}
\end{procedure}

\begin{pseudo}[logreg\_predict dataset\_test.csv weights.csv]
VARIABLES, med, mu, sd, MAISONS, theta $\leftarrow$ lire\_poids(argv[2])\\
lignes $\leftarrow$ lire\_csv(argv[1])\\
\\
pour i dans lignes :\\
\hspace*{2em}si lignes[i][Astronomy] est vide et lignes[i][DADA] est rempli :\\
\hspace*{4em}lignes[i][Astronomy] $\leftarrow$ -100 * lignes[i][DADA]\\
\hspace*{2em}pour f dans VARIABLES :\\
\hspace*{4em}si lignes[i][f] est vide : lignes[i][f] $\leftarrow$ med[f]\\
\\
ouvrir sortie "houses.csv"\\
ecrire "Index,Hogwarts House"\\
\\
pour i dans lignes :\\
\hspace*{2em}x $\leftarrow$ [1.0]\\
\hspace*{2em}pour f dans VARIABLES :\\
\hspace*{4em}ajouter (lignes[i][f] - mu[f]) / sd[f] a x\\
\hspace*{2em}meilleur $\leftarrow$ 0 ; z\_max $\leftarrow$ -infini\\
\hspace*{2em}pour k de 0 a 3 :\\
\hspace*{4em}z $\leftarrow$ 0\\
\hspace*{4em}pour j de 0 a 10 : z $\leftarrow$ z + theta[k][j] * x[j]\\
\hspace*{4em}si z > z\_max : z\_max $\leftarrow$ z ; meilleur $\leftarrow$ k\\
\hspace*{2em}ecrire lignes[i][Index] + "," + MAISONS[meilleur]
\end{pseudo}

\begin{fichier}[houses.csv]
Index,Hogwarts House\\
0,Gryffindor\\
1,Hufflepuff\\
2,Ravenclaw\\
3,Hufflepuff\\
{[}...{]}\\
399,Gryffindor
\end{fichier}

La boucle sur \texttt{k} réalise l'argmax de l'équation~\eqref{eq:ovr-decision}
en une passe, sans stocker les quatre scores. L'initialisation de \texttt{z\_max}
à $-\infty$ garantit qu'une maison est retenue dès le premier tour, y compris
quand les quatre scores sont négatifs.

Les lignes de \texttt{dataset\_test.csv} dont plusieurs notes manquent reçoivent
une maison comme les autres : leurs cases vides ont été remplies par les médianes
d'apprentissage à l'étape~2. Le fichier produit compte $401$ lignes.

\section{Contrôles}

La colonne \texttt{Hogwarts House} de \texttt{dataset\_test.csv} est vide.
L'exactitude se mesure donc sur un ensemble de test découpé dans
\texttt{dataset\_train.csv}, selon le protocole du chapitre~\ref{ch:evaluation}.
Le sujet fixe le seuil d'acceptation à $0{,}98$, mesuré par l'évaluateur sur
\texttt{dataset\_test.csv}.

\begin{table}[htbp]
\centering
\begin{tabularx}{\textwidth}{X r}
\toprule
contrôle & valeur attendue \\
\midrule
gradient analytique contre différence finie centrée & écart relatif $<10^{-6}$ \\
sigmoïde et perte évaluées en $z=\pm1000$ & valeurs finies \\
exactitude d'une prédiction constante sur Hufflepuff & $0{,}331$ \\
exactitude après \rsym{permutation}{permutation} aléatoire des étiquettes & $\approx0{,}34$ \\
exactitude sur le test interne & $\geq0{,}98$ \\
lignes de \texttt{houses.csv} & $401$ \\
\bottomrule
\end{tabularx}
\caption{Valeurs de référence mesurées sur ce jeu de données.}
\label{tab:checks}
\end{table}

Le contrôle par permutation mélange les maisons entre les lignes avant
l'entraînement. Le chapitre~\ref{ch:evaluation} situe le résultat attendu à
$0{,}25$ pour quatre classes équilibrées. Les effectifs valent ici $529$, $443$,
$327$ et $301$ : privé de tout lien entre notes et maison, le modèle converge
vers la prédiction constante sur la classe majoritaire et atteint $0{,}34$. La
référence d'un contrôle par permutation est la performance de la meilleure règle
constante.

\begin{resultat}[Sortie de l'entraînement]
Quatre vecteurs de onze coefficients, dix médianes, dix moyennes, dix
écarts-types, la liste ordonnée des variables et celle des maisons. Ces $74$
nombres et ces deux listes forment l'intégralité de ce que lit la prédiction.
\end{resultat}

\chapter{Notations et lecture des formules}
\label{ch:notations}

\begin{lecture}
Chaque formule met en jeu des objets d'une nature précise et une opération qui
les relie. Les tableaux de ce chapitre donnent pour chaque symbole sa nature, sa
lecture orale et son rôle dans le modèle.
\end{lecture}

Ce chapitre fixe le sens des symboles et la manière de les nommer à l'oral. La
nature d'une notation gouverne les opérations qu'elle admet : un scalaire, un
vecteur, une matrice, une fonction et un ensemble se manipulent selon des règles
propres. La typographie du document la signale — vecteurs en gras, ensembles en
lettres calligraphiques, paramètres grecs en gras lorsqu'ils forment un vecteur.

Le document place l'indice d'observation en bas, selon la convention statistique :
$\x_i$ est le vecteur de l'élève $i$, $x_{ij}$ sa caractéristique $j$ et $c_i$ sa
classe. Certains cours de machine learning écrivent $\x^{(i)}$ pour séparer le
numéro d'exemple d'une puissance ; l'écriture $x_{ij}$ retenue ici s'accorde
directement avec l'usage matriciel et se dicte plus simplement. Les crochets en
exposant, comme $\thv^{[t]}$, portent le numéro d'itération.

\begin{figure}[htbp]
\centering
\resizebox{\textwidth}{!}{%
\begin{tikzpicture}[
  n/.style={draw=accent,fill=accentlight,rounded corners,align=center,
    minimum height=9mm,text width=27mm},
  a/.style={-{Latex},thick,draw=ink},node distance=6mm]
  \node[n] (raw) {$r_{ij}$\\valeur brute};
  \node[n,right=of raw] (x) {$\x_i$\\vecteur préparé};
  \node[n,right=of x] (z) {$z_k$\\score réel};
  \node[n,right=of z] (q) {$q_k$\\sortie binaire};
  \node[n,right=of q] (c) {$\widehat c$\\classe prédite};
  \draw[a] (raw)--node[above,scale=.78]{standardiser}(x);
  \draw[a] (x)--node[above,scale=.78]{$\thv_k^\top\xt$}(z);
  \draw[a] (z)--node[above,scale=.78]{$\sigma$}(q);
  \draw[a] (q)--node[above,scale=.78]{$\Argmax_k$}(c);
\end{tikzpicture}}
\caption{Nature des objets successifs et transformation qui mène de l'un à
l'autre : une valeur brute, un vecteur préparé, un score réel, une sortie
logistique, une classe.}
\label{fig:notation-chain}
\end{figure}

\section{Indices, effectifs et ensembles}

\begin{table}[htbp]
\centering
\begin{tabularx}{\textwidth}{>{$}l<{$} >{\raggedright\arraybackslash}p{3.1cm} X}
\toprule
\text{symbole} & \text{lecture orale} & \text{signification exacte} \\
\midrule
i & « indice i » lorsqu'il est attaché à un symbole & indice d'une observation : ici, un élève. Il varie de $1$ à $m$. \\
j & « indice j » lorsqu'il est attaché à un symbole & indice d'une caractéristique : ici, un cours ou une note. Il varie de $1$ à $n$. \\
k & « indice k » lorsqu'il est attaché à un symbole & indice d'une classe : ici, une maison. Il varie de $1$ à $K$. \\
t & « itération t » & indice d'une itération de l'algorithme d'optimisation. \\
m & « m » & nombre total d'observations dans l'ensemble considéré. \\
n & « n » & nombre de caractéristiques utilisées par le modèle, avant ajout du biais. \\
K & « K » & nombre de classes ; dans ce problème, $K=4$. \\
\mathcal C & « l'ensemble des classes C » & ensemble des classes possibles, et non une classe particulière. On précise « C calligraphique » seulement si l'on doit dicter la graphie. \\
\R & « R » ou « l'ensemble des réels » & ensemble dans lequel prennent leurs valeurs les scores, paramètres et caractéristiques numériques. \\
\mathcal I_{\mathrm{tr}} & « I indice apprentissage » & ensemble des indices appartenant à l'échantillon d'apprentissage. \\
\mathcal I_{\mathrm{te}} & « I indice test » & ensemble des indices appartenant à l'échantillon de test. \\
\bottomrule
\end{tabularx}
\caption{Indices, tailles et ensembles. Un même indice conserve le même rôle dans
tout le document.}
\end{table}

L'écriture $i\in\mathcal I_{\mathrm{tr}}$ se lit « $i$ appartient à l'ensemble
des indices d'apprentissage ». L'écriture $k\in\{1,\ldots,K\}$ se lit « $k$ varie
de un à $K$ ».

\section{Une observation et ses caractéristiques}

\begin{table}[htbp]
\centering
\begin{tabularx}{\textwidth}{>{$}l<{$} >{\raggedright\arraybackslash}p{3.4cm} X}
\toprule
\text{symbole} & \text{lecture orale} & \text{signification exacte} \\
\midrule
r_{ij} & « r indice i j » & valeur brute de la caractéristique $j$ pour l'élève $i$, avant standardisation. Le premier indice repère la ligne, le second la variable. \\
\mu_j & « mu indice j » & moyenne de la caractéristique $j$, estimée uniquement sur l'apprentissage. \\
s_j & « s indice j » & écart-type de la caractéristique $j$, estimé uniquement sur l'apprentissage. \\
x_{ij} & « x indice i j » & caractéristique $j$ standardisée de l'observation $i$ ; c'est un scalaire. \\
\x_i & « x indice i » & vecteur colonne des $n$ caractéristiques standardisées de l'observation $i$ ; $\x_i\in\R^n$. Le caractère gras est généralement implicite à l'oral. \\
\xt_i & « x tilde indice i » ou « x i augmenté » & vecteur $\x_i$ précédé d'une coordonnée égale à $1$ pour intégrer le biais ; $\xt_i\in\R^{n+1}$. \\
\X & « la matrice X » ou « la matrice de design » & matrice qui empile les vecteurs augmentés en lignes ; $\X\in\R^{m\times(n+1)}$. \\
X_i & « la variable aléatoire X indice i » & vecteur aléatoire représentant les caractéristiques de l'observation $i$ avant leur réalisation. Il se distingue de la valeur observée $\x_i$ et de la matrice $\X$. \\
\bottomrule
\end{tabularx}
\caption{Les minuscules désignent une observation ou une composante ; la
majuscule grasse $\X$ réunit toutes les observations sous forme matricielle.}
\end{table}

La formule
\[
  x_{ij}=\frac{r_{ij}-\mu_j}{s_j}
\]
se lit : « x indice i j égale r indice i j moins mu indice j, sur s indice j ».
Elle exprime la valeur de la note en nombre d'écarts-types au-dessus ou au-dessous
de la moyenne d'apprentissage.

\section{Classes observées, cibles binaires et prédictions}

\begin{table}[htbp]
\centering
\begin{tabularx}{\textwidth}{>{$}l<{$} >{\raggedright\arraybackslash}p{3.5cm} X}
\toprule
\text{symbole} & \text{lecture orale} & \text{signification exacte} \\
\midrule
C_i & « la variable aléatoire C indice i » & maison de l'élève $i$ avant observation ; $C_i$ prend ses valeurs dans $\mathcal C$. \\
c_i & « c indice i » & maison effectivement observée pour l'élève $i$ ; $c_i\in\mathcal C$. La minuscule désigne ici une réalisation de $C_i$. \\
y_{ik} & « y indice i k » & cible binaire du classifieur $k$ pour l'observation $i$ : $1$ si $c_i=k$, sinon $0$. \\
\y_k & « le vecteur y indice k » & vecteur colonne réunissant les $m$ cibles du problème « classe $k$ contre le reste » ; $\y_k\in\{0,1\}^m$. \\
\widehat y_i & « y chapeau indice i » & classe binaire prédite pour l'observation $i$, égale à $0$ ou $1$. Le chapeau marque une valeur estimée. \\
\widehat c(\x) & « c chapeau évalué en x » ou « la classe prédite en x » & règle de classification multiclasses évaluée pour les caractéristiques $\x$. \\
\ind[A] & « l'indicatrice de l'événement A » & vaut $1$ lorsque la proposition $A$ est vraie et $0$ sinon. \\
\bottomrule
\end{tabularx}
\caption{La classe observée $c_i$ prend ses valeurs dans $\mathcal C$ ; la cible
binaire $y_{ik}$, construite pour le sous-problème OvR numéro $k$, vaut $0$ ou
$1$.}
\end{table}

L'écriture $y_{ik}=\ind[c_i=k]$ se lit : « y indice i k est l'indicatrice de
l'événement c indice i égale k » ; en contexte, on dira plus naturellement :
« la cible de l'observation i pour la classe k vaut un si sa classe est k ».

\section{Paramètres, score et probabilité binaire}

\begin{table}[htbp]
\centering
\begin{tabularx}{\textwidth}{>{$}l<{$} >{\raggedright\arraybackslash}p{3.7cm} X}
\toprule
\text{symbole} & \text{lecture orale} & \text{signification exacte} \\
\midrule
\theta_j & « thêta indice j » & coefficient scalaire associé à la caractéristique $j$. Son signe indique le sens de son effet sur les log-cotes, toutes choses égales par ailleurs. \\
\theta_0 & « thêta zéro » & biais ou constante du modèle ; il multiplie la coordonnée ajoutée $x_0=1$. \\
\thv & « thêta » ou « le vecteur des paramètres thêta » & vecteur colonne de tous les paramètres $(\theta_0,\ldots,\theta_n)^\top\in\R^{n+1}$. On ne dit « vecteur » que lorsque sa nature doit être soulignée. \\
\thv_k & « thêta indice k » & vecteur des paramètres du classifieur binaire « classe $k$ contre le reste ». \\
z_{\thv}(\x) & « le score z, paramétré par thêta, évalué en x » ; souvent « le score en x » & score linéaire réel $\thv^\top\xt$, à valeurs dans $\R$ tout entier. \\
\sigma & « sigma » & fonction logistique $z\mapsto(1+e^{-z})^{-1}$, qui transforme un score réel en nombre strictement compris entre $0$ et $1$. \\
p_{\thv}(\x) & « p paramétré par thêta, évalué en x » ; souvent « la probabilité prédite en x » & probabilité conditionnelle modélisée de la classe binaire positive : $\Pp(Y=1\mid X=\x)$. \\
q_k(\x) & « q indice k évalué en x » & sortie du classifieur binaire $k$, égale à $\sigma(\thv_k^\top\xt)$. La somme des $K$ valeurs OvR prend une valeur libre. \\
\Theta & « thêta majuscule » ou « la matrice des paramètres » & matrice regroupant les $K$ vecteurs de paramètres, un par classifieur OvR. \\
\bottomrule
\end{tabularx}
\caption{Le modèle produit successivement trois objets : le score $z$ dans $\R$,
la probabilité binaire $p$ dans $]0,1[$, la décision $\widehat y$ dans
$\{0,1\}$.}
\end{table}

Le produit $\thv^\top\xt$ se lit « thêta transposé x tilde », ou « le produit
scalaire de thêta et x tilde ». Le symbole
$\top$ indique la transposition : il transforme ici le vecteur colonne $\thv$ en
vecteur ligne afin que le produit scalaire soit défini. Développé, ce produit vaut
\[
  \thv^\top\xt=\theta_0+\theta_1x_1+\cdots+\theta_nx_n.
\]

\section{Vraisemblance, perte et optimisation}

\begin{table}[htbp]
\centering
\begin{tabularx}{\textwidth}{>{$}l<{$} >{\raggedright\arraybackslash}p{3.7cm} X}
\toprule
\text{symbole} & \text{lecture orale} & \text{signification exacte} \\
\midrule
L(\thv) & « la vraisemblance évaluée en thêta » & vraisemblance : probabilité des étiquettes observées, considérée comme fonction des paramètres. \\
\log L(\thv) & « la log-vraisemblance en thêta » & logarithme de la vraisemblance ; il remplace les produits par des sommes sans changer le maximiseur. \\
\ell_i(\thv) & « ell indice i évaluée en thêta » ou « la perte de l'observation i » & opposé de la contribution de l'observation $i$ à la log-vraisemblance. \\
J(\thv) & « J évaluée en thêta » ou « le coût en thêta » & fonction objectif : moyenne des pertes sur l'échantillon d'apprentissage. \\
\nabla J(\thv) & « le gradient de J en thêta » & vecteur des $n+1$ dérivées partielles de $J$ ; il pointe dans la direction locale de plus forte augmentation. « Nabla J » sert à dicter la formule. \\
\nabla^2J(\thv) & « la Hessienne de J en thêta » & matrice $(n+1)\times(n+1)$ des dérivées secondes ; elle décrit la courbure locale de la fonction objectif. « Nabla carré J » nomme la notation, « Hessienne » nomme l'objet. \\
W & « la matrice W » & matrice diagonale dont le terme $i$ vaut $p_i(1-p_i)$ ; elle intervient dans la Hessienne. \\
\alpha & « alpha » & taux d'apprentissage, c'est-à-dire longueur multiplicative du pas de descente de gradient. \\
\lambda & « lambda » & intensité de la pénalisation quadratique ; $\lambda=0$ signifie absence de régularisation. \\
\eps & « epsilon » & tolérance numérique utilisée dans le critère d'arrêt. \\
\bottomrule
\end{tabularx}
\caption{Objets employés pour estimer les paramètres. La vraisemblance est
maximisée ; son opposé logarithmique $J$ est minimisé.}
\end{table}

La formule
\[
  \nabla J(\thv)=\frac1m\X^\top\bigl(\sigma(\X\thv)-\y\bigr)
\]
se lit : « le gradient de J en thêta vaut un sur m, X transposé, multiplié par
le vecteur sigma de X thêta moins y ». Le vecteur entre parenthèses contient les
écarts entre probabilités prédites et étiquettes observées.

La mise à jour
\[
  \thv^{[t+1]}=\thv^{[t]}-\alpha\nabla J(\thv^{[t]})
\]
se lit : « à l'itération t plus un, thêta vaut sa valeur à l'itération t moins
alpha fois le gradient évalué en cette valeur ». Les crochets $[t]$ repèrent le
numéro d'itération.

\section{Probabilités et opérateurs}

\begin{table}[htbp]
\centering
\begin{tabularx}{\textwidth}{>{$}l<{$} >{\raggedright\arraybackslash}p{3.9cm} X}
\toprule
\text{symbole} & \text{lecture orale} & \text{signification exacte} \\
\midrule
\Pp(A) & « probabilité de l'événement A » & probabilité de l'événement $A$. \\
\Pp(A\mid B) & « probabilité de A conditionnellement à B » ; couramment « A sachant B » & probabilité conditionnelle de $A$ lorsque l'information $B$ est connue. \\
Y\mid X=\x & « la loi de Y conditionnellement à X égal x » & loi conditionnelle de la réponse pour une observation dont les caractéristiques valent $\x$. \\
\sum_{i=1}^m & « somme sur i, de un à m » & addition de l'expression située à droite pour toutes les observations. \\
\prod_{i=1}^m & « produit sur i, de un à m » & multiplication de l'expression située à droite pour toutes les observations. \\
\Argmax_k a_k & « l'argument du maximum sur k de a indice k » & indice $k$ pour lequel $a_k$ atteint son maximum ; le résultat est cet indice, quand $\max_k a_k$ désigne la valeur atteinte. Dans l'usage, « arg max » est également courant. \\
\|v\|_2 & « norme euclidienne de v » ou « norme deux de v » & longueur euclidienne du vecteur $v$. \\
\operatorname{diag}(a_i) & « la matrice diagonale de termes a indice i » & matrice dont la diagonale contient les $a_i$ et dont les autres termes sont nuls. \\
\partial J/\partial\theta_j & « dérivée partielle de J par rapport à thêta indice j » & variation locale de $J$ lorsque seul $\theta_j$ varie, les autres paramètres étant maintenus fixes. \\
\bottomrule
\end{tabularx}
\caption{Opérateurs qui apparaissent dans les dérivations et les règles de
décision.}
\end{table}

Enfin,
\[
  \widehat c(\x)=\Argmax_{k\in\{1,\ldots,K\}}\thv_k^\top\xt
\]
se lit : « la classe prédite en x est l'argument qui maximise, sur les classes k,
le score thêta indice k transposé x tilde ». Cette formule renvoie le nom de la
classe gagnante après conversion de l'indice vers le libellé correspondant.

\chapter{Guide de lecture et ressources ouvertes}
\label{ch:reading-guide}

Une bibliographie brute ne dit pas par où commencer. Les ressources ci-dessous
sont donc ordonnées par objectif et par niveau mathématique. Tous les intitulés
en couleur sont des liens cliquables vers une version ouverte ou vers la page
officielle de la ressource.

\section{Premiers repères statistiques}

\begin{description}[style=nextline,leftmargin=0pt]
  \item[\href{https://openintrostat.github.io/ims/}{Introduction to Modern
  Statistics --- chapitre 9}] Ressource libre de Çetinkaya-Rundel et Hardin
  \cite{openintro2024}. Le chapitre sur la régression logistique part de données,
  explique les cotes et l'interprétation des coefficients avant de développer
  l'inférence. C'est le meilleur point d'entrée de cette liste si les probabilités
  conditionnelles et les log-cotes ne sont pas encore familières.

  \item[\href{https://www.statlearning.com/}{An Introduction to Statistical
  Learning}] Ouvrage de James, Witten, Hastie, Tibshirani et Taylor
  \cite{james2023}, disponible gratuitement en versions R et Python. Le traitement
  est moins technique que \emph{The Elements of Statistical Learning}, avec des
  exemples, des figures et des laboratoires. Lire le chapitre Classification,
  puis les chapitres Resampling Methods et Linear Model Selection and
  Regularization pour comprendre validation croisée et pénalisation.
\end{description}

\section{Probabilités et calcul}

\begin{description}[style=nextline,leftmargin=0pt]
  \item[\href{https://web.stanford.edu/class/archive/cs/cs109/cs109.1264/lectures/20-LogisticRegression/}{Stanford CS109 --- Logistic Regression}]
  Cours court accompagné de diapositives et de code \cite{stanfordcs109}. Il est
  utile pour revoir maximum de vraisemblance et régression logistique dans une
  progression de probabilités appliquées, sans entrer immédiatement dans un
  traité de statistique mathématique.

  \item[\href{https://see.stanford.edu/Course/CS229}{Stanford Engineering
  Everywhere --- CS229}] Le cours historique fournit vidéos, transcriptions,
  exercices et notes. Les notes 1 couvrent classification et régression
  logistique ; les notes 5 couvrent régularisation et sélection de modèle
  \cite{ng2018cs229}. La densité mathématique est supérieure aux deux ressources
  précédentes, mais les dérivations sont directement utilisables.

  \item[\href{https://online.stat.psu.edu/stat504/}{Penn State STAT 504 ---
  Analysis of Discrete Data}] Notes de cours ouvertes du département de
  statistique de Penn State \cite{psustat504}. Les leçons 6 à 8 traitent
  successivement régression logistique binaire, diagnostics et régression
  multinomiale. Cette ressource est particulièrement utile pour approfondir
  rapports de cotes, qualité d'ajustement et interprétation statistique.
\end{description}

\section{Approfondir le formalisme}

\begin{description}[style=nextline,leftmargin=0pt]
  \item[\href{https://hastie.su.domains/ElemStatLearn/}{The Elements of
  Statistical Learning}] Référence de Hastie, Tibshirani et Friedman
  \cite{hastie2009}. Le chapitre 4 contient les équations citées dans ce document.
  Il est concis et exigeant : mieux vaut l'utiliser après une première lecture
  d'ISL ou d'OpenIntro, puis vérifier une à une vraisemblance, gradient, Hessienne
  et régularisation.

  \item[\href{https://probml.github.io/pml-book/book1.html}{Probabilistic Machine
  Learning: An Introduction}] Ouvrage de Murphy \cite{murphy2022}. Il replace la
  régression logistique dans un cadre probabiliste plus large et développe
  optimisation, calibration et modèles multiclasses. À consulter lorsque les
  objets du présent cours sont déjà maîtrisés séparément.

  \item[\href{https://www.microsoft.com/en-us/research/publication/pattern-recognition-machine-learning/}{Pattern
  Recognition and Machine Learning}] Ouvrage de Bishop \cite{bishop2006}. Sa
  section 4.3 traite les modèles linéaires de classification depuis un point de
  vue probabiliste, et développe la méthode de Newton--Raphson comme alternative
  à la descente de gradient. C'est la référence à consulter pour dépasser la
  descente du premier ordre employée ici.

  \item[\href{https://web.stanford.edu/~boyd/cvxbook/}{Convex Optimization}]
  Ouvrage ouvert de Boyd et Vandenberghe \cite{boyd2004}. Les chapitres sur la
  convexité et les méthodes de descente donnent le cadre général derrière la
  Hessienne positive et le choix d'une direction de descente. Cette référence
  traite l'optimisation mathématique générale ; son étude suppose donc déjà le
  contexte statistique de la classification.
\end{description}

\begin{resultat}[Parcours conseillé]
Pour une première étude : OpenIntro, puis ISL, puis les notes CS229. Pour vérifier
le formalisme du présent document : ESL chapitre 4 et Boyd chapitre 9. Pour aller
au-delà d'OvR vers une modélisation probabiliste multiclasses : Penn State leçon
8, puis Murphy. Cette progression évite de demander à une même source d'être à la
fois introduction, manuel de calcul et référence théorique.
\end{resultat}

\chapter*{Conclusion}
\addcontentsline{toc}{chapter}{Conclusion}

La régression logistique repose sur une chaîne cohérente : modèle linéaire des
log-cotes, loi de Bernoulli conditionnelle, maximum de vraisemblance, log-perte
convexe et optimisation numérique contrôlée. OvR réutilise ce problème binaire
$K$ fois et fournit une règle multiclasses par maximum des scores. Sa simplicité
a une contrepartie précise : la somme des sorties indépendantes prend une valeur
libre, quand une distribution catégorielle vaut $1$. La validité du résultat
tient enfin autant au protocole expérimental --- séparation stricte du test,
diagnostics par classe, archivage --- qu'à la correction des équations.

\begin{resultat}[Enchaînement des opérations]
Le modèle reçoit des caractéristiques préparées, calcule quatre combinaisons
linéaires et retient le score maximal. Son apprentissage choisit chaque vecteur
de paramètres de façon à maximiser la vraisemblance de la cible binaire associée.
Sa validation porte sur trois questions séparées : la convergence numérique de
l'optimisation, la généralisation hors échantillon, et les erreurs propres à
chaque maison. Chacune se mesure par ses propres quantités.
\end{resultat}

\appendix
\chapter{Dérivations mathématiques}
\label{app:derivations}

Cette annexe rassemble les justifications algébriques des formules employées dans
le corps du document. Les chapitres principaux en gardent les expressions et leur
interprétation opérationnelle.

\section{Inversion du logit}
\label{app:logit-inversion}

En posant $p=\Pp(Y=1\mid X=\x)$ et $z=\thv^\top\xt$,
\begin{align*}
  \log\frac{p}{1-p}=z
  &\iff \frac{p}{1-p}=e^z
  \iff p=e^z(1-p)\\
  &\iff p(1+e^z)=e^z
  \iff p=\frac{e^z}{1+e^z}=\frac1{1+e^{-z}}.
\end{align*}

\section{Gradient de la log-perte}
\label{app:gradient-derivation}

Pour une observation,
$\ell_i=-y_i\log p_i-(1-y_i)\log(1-p_i)$ et
$\partial p_i/\partial z_i=p_i(1-p_i)$. La règle de chaîne donne
\begin{align*}
  \frac{\partial\ell_i}{\partial z_i}
  &=-\frac{y_i}{p_i}p_i(1-p_i)
    +\frac{1-y_i}{1-p_i}p_i(1-p_i)\\
  &=p_i-y_i.
\end{align*}
Comme $\nabla_{\thv}z_i=\xt_i$, la moyenne des contributions vaut
\[
  \nabla J(\thv)=\frac1m\sum_{i=1}^m(p_i-y_i)\xt_i
  =\frac1m\X^\top\bigl(\sigma(\X\thv)-\y\bigr).
\]

\section{Convexité de la log-perte}
\label{app:convexity-proof}

Avec $W=\operatorname{diag}(p_i(1-p_i))$,
\[
  \nabla^2J(\thv)=\frac1m\X^\top W\X.
\]
Pour tout $v\in\R^{n+1}$,
\[
  v^\top\nabla^2J(\thv)v
  =\frac1m(\X v)^\top W(\X v)
  =\frac1m\sum_{i=1}^m p_i(1-p_i)(\xt_i^\top v)^2\geq0.
\]
La Hessienne est donc positive semi-définie et $J$ est convexe.

\section{Direction de descente}
\label{app:descent-direction}

Pour un déplacement suffisamment petit $\Delta\thv$,
\[
  J(\thv+\Delta\thv)
  =J(\thv)+\nabla J(\thv)^\top\Delta\thv
  +o(\|\Delta\thv\|).
\]
Avec $\Delta\thv=-\alpha\nabla J(\thv)$,
\[
  J(\thv+\Delta\thv)
  =J(\thv)-\alpha\|\nabla J(\thv)\|_2^2
  +o(\alpha\|\nabla J(\thv)\|).
\]
Le terme principal est négatif pour $\alpha>0$ lorsque le gradient est non nul.


\printbibliography[title={Références citées},notcategory=reading]
\nocite{bishop2006,ng2018cs229,murphy2022,james2023,
  openintro2024,stanfordcs109,psustat504}
\printbibliography[title={Lectures complémentaires},category=reading]
\end{document}
